\documentclass[11pt,a4paper,oneside,openany]{book}
\usepackage[T1]{fontenc}
\usepackage[utf8]{inputenc}
\usepackage{libertinus}
\usepackage{microtype}
\usepackage[a4paper,margin=27mm,headheight=15pt]{geometry}
\usepackage{amsmath,amssymb,amsthm,mathtools,bm,mathrsfs}
\usepackage{booktabs,longtable,array,tabularx}
\usepackage{enumitem,xcolor}
\usepackage[most]{tcolorbox}
\usepackage{fancyhdr,xurl,hyperref}
\usepackage[nameinlink,noabbrev,capitalise]{cleveref}
\definecolor{deepolive}{RGB}{67,79,45}
\definecolor{softolive}{RGB}{236,239,226}
\definecolor{darkteal}{RGB}{31,83,84}
\definecolor{warmgray}{RGB}{92,87,78}
\definecolor{softcoral}{RGB}{246,225,219}
\definecolor{linkblue}{RGB}{31,76,122}
\hypersetup{colorlinks=true,linkcolor=deepolive,citecolor=darkteal,urlcolor=linkblue,
 pdftitle={q-Binomial Coefficient Calculus: An Extension of Combinatorial Coefficient Calculus},
 pdfauthor={A proof-oriented mathematical extension},
 pdfsubject={Gaussian coefficients, Bell polynomials, q-Stirling numbers, parameter jets, and inversion}}
\allowdisplaybreaks[3]
\raggedbottom
\emergencystretch=2em
\setlist{itemsep=2pt,topsep=4pt,parsep=1pt}
\setcounter{tocdepth}{1}
\setcounter{secnumdepth}{2}
\pagestyle{fancy}
\fancyhf{}
\fancyhead[L]{\small\nouppercase{\leftmark}}
\fancyhead[R]{\small\thepage}
\renewcommand{\headrulewidth}{0.4pt}
\renewcommand{\chaptermark}[1]{\markboth{\thechapter.\ #1}{}}
\fancypagestyle{plain}{\fancyhf{}\fancyfoot[C]{\thepage}\renewcommand{\headrulewidth}{0pt}}
\newtheorem{theorem}{Theorem}[chapter]
\newtheorem{lemma}[theorem]{Lemma}
\newtheorem{proposition}[theorem]{Proposition}
\newtheorem{corollary}[theorem]{Corollary}
\theoremstyle{definition}
\newtheorem{definition}[theorem]{Definition}
\newtheorem{example}[theorem]{Example}
\theoremstyle{remark}
\newtheorem{remark}[theorem]{Remark}
\newtcolorbox{masterbox}[1][]{enhanced,breakable,colback=softolive,colframe=deepolive,
 boxrule=0.7pt,arc=1mm,left=2mm,right=2mm,top=1.5mm,bottom=1.5mm,
 title={Master principle},fonttitle=\bfseries,#1}
\newtcolorbox{auditbox}[1][]{enhanced,breakable,colback=softcoral,colframe=warmgray,
 boxrule=0.6pt,arc=1mm,left=2mm,right=2mm,top=1.5mm,bottom=1.5mm,
 title={Normalization and domain},fonttitle=\bfseries,#1}
\newcommand{\N}{\mathbb N}
\newcommand{\Z}{\mathbb Z}
\newcommand{\Q}{\mathbb Q}
\newcommand{\C}{\mathbb C}
\newcommand{\R}{\mathbb R}
\newcommand{\F}{\mathbb F}
\newcommand{\GB}[3][q]{\begin{bmatrix}#2\\#3\end{bmatrix}_{#1}}
\newcommand{\qi}[2][q]{[#2]_{#1}}
\newcommand{\qf}[2][q]{[#2]_{#1}!}
\newcommand{\Y}[1]{\mathcal B_{#1}}
\newcommand{\B}[2]{\mathcal B_{#1,#2}}
\newcommand{\qB}[2]{\mathfrak B^{[q]}_{#1,#2}}
\newcommand{\coeff}[2]{[#1^{#2}]}
\newcommand{\fall}[2]{#1^{\underline{#2}}}
\newcommand{\rising}[2]{#1^{\overline{#2}}}
\DeclareMathOperator{\Li}{Li}
\DeclareMathOperator{\Var}{Var}
\DeclareMathOperator{\ord}{ord}
\DeclareMathOperator{\Res}{Res}
\newcommand{\dd}{\,\mathrm d}
\newcommand{\E}{\mathbb E}
\newcommand{\qbname}{\texorpdfstring{$q$}{q}}
\newcommand{\logq}{\vartheta}
\title{\Huge\bfseries \(q\)-Binomial Coefficient Calculus\\[4mm]
 \Large Gaussian Kernels, Bell Polynomials,\\
 Basis Changes, Parameter Jets, and Inversion\\[7mm]
 \large An extension of \emph{Combinatorial Coefficient Calculus}}
\author{A self-contained development with detailed proofs}
\date{September 4, 2026}
\begin{document}
\frontmatter
\hypersetup{pageanchor=false}
\maketitle
\hypersetup{pageanchor=true}
\chapter*{Abstract}
\addcontentsline{toc}{chapter}{Abstract}
This article extends the basis-change and coefficient-extraction framework of
\emph{Combinatorial Coefficient Calculus} to Gaussian, or \(q\)-binomial,
coefficients. The extension is organized around mathematical operations, rather
than a list of disconnected identities. Finite products produce Gaussian kernels;
logarithms of those products produce power sums; Bell polynomials recover the
coefficients; Newton interpolation produces \(q\)-Stirling arrays; and ordinary
formal composition and Lagrange inversion transport the resulting formulae to
implicit series.

The development includes Gaussian convolution and inversion, finite-field
linearization, terminating \(q\)-Chu--Vandermonde sums, principal specializations
of symmetric functions, a master coefficient formula for arbitrary powers of
finite geometric products, partial-fraction proofs of divisor-type identities,
Jackson differentiation and geometric interpolation, two explicitly distinguished
\(q\)-Stirling normalizations, and a weighted Bell--Riordan calculus. It also
provides recurrence-free coefficients of Gaussian polynomials through divisor
sums, all-order expansions at \(q=1\), and a uniform deflation procedure giving
all Taylor coefficients at every root of unity.

The last analytic application constructs Gaussian extrapolation filters. Their
complete error response and exact noise norm are computed, and repeated spectral
roots extend the construction to polynomial--logarithmic errors. Every theorem
used in the development is proved or reduced explicitly to a proved result.
Classical identities are identified as such; the synthesis and further
consequences are not presented as claims of historical priority. A companion
standard-library Python program performs exact, deterministic normalization
checks. These checks supplement, rather than replace, the proofs.

\chapter*{Relation to the source article}
\addcontentsline{toc}{chapter}{Relation to the source article}
The source article \cite{source} develops two central mechanisms: changing to a
basis adapted to an operator, and using Bell polynomials to recover coefficients
from component data. It also separates formal statements from analytic statements
and keeps Bernoulli numbers distinct from Bell polynomials. This extension retains
that organization and those distinctions.

Its first point of attachment is the source's Newton expansion: the additive
nodes \(0,1,2,\ldots\) are replaced either by geometric nodes \(a,aq,aq^2,\ldots\)
or by \(q\)-integer nodes \([0]_q,[1]_q,[2]_q,\ldots\). These are related but not
identical constructions. Its second point of attachment is the exponential
formula: a finite geometric alphabet has power sums \([N]_{q^r}\), so the same
ordinary Bell polynomials already used in the source immediately encode
\(q\)-binomial identities. Its third point of attachment is ordinary composition
and reversion. A change to \(q\)-factorial coefficient normalization changes the
matrices, not the meaning of ordinary substitution.

This is a standalone mathematical extension, not a replacement of the canonical
article and not a claim that its additional results have been formalized in Lean.
All notation required for compilation is defined in this file; the repository's
shared notation file is not needed. The typography follows the source's
Libertinus, olive, and coral conventions.

\begin{auditbox}[title={What ``an extension'' means here}]
The finite \(q\)-binomial theorem, Gaussian inversion, \(q\)-Chu--Vandermonde,
Carlitz-type arrays, divisor identities, and \(q\)-Lucas are established
mathematics. They are rederived to connect their hypotheses and normalizations.
The main additional value is the common coefficient machinery, including
explicit all-order root-of-unity jets and full geometric-filter error formulae.
No exhaustive novelty claim is made.
\end{auditbox}
\begingroup
\small
\tableofcontents
\endgroup

\mainmatter
\chapter{The coefficient-calculus setting}
\label{qcc:setting}
\section{Rings, indices, and specialization}
Unless stated otherwise, \(q\) is an indeterminate and calculations with
quotients take place over \(\Q(q)\), extended by the finitely many auxiliary
parameters that occur. A formal series in \(z\) is an element of \(K[[z]]\) for a
commutative \(\Q\)-algebra \(K\). Composition \(A(B(z))\) is defined when
\(B(0)=0\); inversion under multiplication requires an invertible constant term.
These statements have no convergence content.

All integer indices are nonnegative unless explicitly qualified. Empty products
are one, sums over empty ranges are zero, and \(0^0=1\) in coefficient formulae.
Gaussian coefficients outside \(0\le k\le n\) are zero. Define
\begin{equation}\label{qcc:notation}
 (a;q)_N=\prod_{j=0}^{N-1}(1-aq^j),\qquad
 \qi n=\sum_{j=0}^{n-1}q^j,\qquad
 \qf n=\prod_{j=1}^n\qi j.
\end{equation}
In particular, \([0]_q=0\) and \([0]_q!=1\). The polynomial interpretation of
\([n]_q\) makes \([n]_1=n\) immediate. The quotient
\((1-q^n)/(1-q)\) is a useful representation, not the definition at \(q=1\).
Our \(q\)-shifted factorial convention agrees with \cite{dlmf172}.

When a rational representation equals a polynomial, specialization means
specializing that polynomial, after cancellation. In particular, the expression
\((q;q)_n/((q;q)_k(q;q)_{n-k})\) must not be evaluated as an uncancelled
\(0/0\) at a root of unity. Later we give explicit polynomial and local-series
methods that do not make this mistake.

\section{Bell and Bernoulli normalizations}
The complete and partial exponential Bell polynomials are characterized by
\begin{align}
 \exp\left(\sum_{r\ge1}x_r\frac{z^r}{r!}\right)
 &=\sum_{n\ge0}\Y n(x_1,\ldots,x_n)\frac{z^n}{n!},\label{qcc:bell-def}\\
 \frac1{k!}\left(\sum_{r\ge1}x_r\frac{z^r}{r!}\right)^k
 &=\sum_{n\ge k}\B nk(x_1,\ldots)\frac{z^n}{n!}.\label{qcc:partial-def}
\end{align}
Thus \(\Y0=1\), and
\begin{equation}\label{qcc:bell-finite}
 \Y n(x_1,\ldots,x_n)
 =n!\!\sum_{\substack{m_1,\ldots,m_n\ge0\\\sum r m_r=n}}
 \prod_{r=1}^n\frac1{m_r!}\left(\frac{x_r}{r!}\right)^{m_r}.
\end{equation}
For \(\B nk\), impose also \(\sum m_r=k\). Expanding each exponential in
\eqref{qcc:bell-def} proves \eqref{qcc:bell-finite}. In particular, a formula in
Bell polynomials is a finite-sum formula, not a hidden recurrence.

Bernoulli numbers are denoted \(\beta_r\), with
\begin{equation}\label{qcc:bernoulli-def}
 \frac{t}{e^t-1}=\sum_{r\ge0}\beta_r\frac{t^r}{r!},
 \qquad \beta_1=-\frac12.
\end{equation}
We use \(s(n,k)\) for signed ordinary Stirling numbers of the first kind and
\(S(n,k)\) for ordinary Stirling numbers of the second kind. Their \(q\)-versions
will be defined independently, with the normalization stated in the definition.

\section{Three distinct operations}
There are three operations that should not be conflated. One may specialize the
variables of a symmetric polynomial to a geometric progression. One may replace
an ordinary derivative by the Jackson derivative. Or one may merely record
ordinary series coefficients using \([n]_q!\) rather than \(n!\). All three lead
to Gaussian coefficients, but their product and composition rules are not
interchangeable.

\begin{masterbox}
The robust route is
\[
 \text{finite alphabet}\ \longrightarrow\ \text{power sums}
 \ \longrightarrow\ \text{formal logarithm}
 \ \longrightarrow\ \text{Bell coefficients}.
\]
Operator-adapted interpolation and formal reversion can then be attached to this
route without guessing which ordinary factorials should acquire a subscript \(q\).
\end{masterbox}

\chapter{Gaussian polynomials and finite products}
\label{qcc:gaussian}
\section{A division-free definition}
\begin{definition}
For \(0\le k\le n\), set
\begin{equation}\label{qcc:subset}
 \GB nk=\sum_{0\le i_1<\cdots<i_k<n}
 q^{i_1+\cdots+i_k-\binom{k}{2}}.
\end{equation}
The \(k=0\) summand is the empty subset and contributes one.
\end{definition}
Every exponent is nonnegative, because \(i_j\ge j-1\). Consequently this is a
polynomial in \(\Z_{\ge0}[q]\), valid at every complex \(q\).

\begin{theorem}[Pascal, factorial, and symmetry formulae]\label{qcc:pascal}
For \(n\ge1\), the Gaussian polynomials satisfy the first two recurrences
below; the remaining identities hold for \(0\le k\le n\):
\begin{align}
 \GB nk&=\GB{n-1}k+q^{n-k}\GB{n-1}{k-1},\label{qcc:pascal1}\\
 \GB nk&=q^k\GB{n-1}k+\GB{n-1}{k-1},\label{qcc:pascal2}\\
 \GB nk&=\frac{(q;q)_n}{(q;q)_k(q;q)_{n-k}}
       =\frac{\qf n}{\qf k\qf{n-k}},\label{qcc:quotient}\\
 \GB nk&=\GB n{n-k},\qquad
 \GB[q^{-1}]nk=q^{-k(n-k)}\GB nk.\label{qcc:reciprocity}
\end{align}
The quotient identities hold in \(\Q(q)\). For \(0\le k\le n\), the polynomial
has degree \(k(n-k)\), constant coefficient one, and leading coefficient one.
\end{theorem}
\begin{proof}
Separate the subsets in \eqref{qcc:subset} according as they contain \(n-1\).
Removing that element decreases the exponent by \(n-k\), proving
\eqref{qcc:pascal1}. The factorial quotient has the same boundary values and
satisfies the same recurrence: after putting the two terms over a common
denominator the remaining numerator is
\((1-q^{n-k})+q^{n-k}(1-q^k)=1-q^n\).
Induction therefore proves \eqref{qcc:quotient}. Its symmetry gives
\eqref{qcc:pascal2} by replacing \(k\) with \(n-k\).

The smallest subset \(\{0,\ldots,k-1\}\) has exponent zero, and the largest
\(\{n-k,\ldots,n-1\}\) has exponent \(k(n-k)\). Each is unique.
Reflecting indices by \(i_j\mapsto n-1-i_{k+1-j}\) replaces an exponent \(d\)
by \(k(n-k)-d\), proving reciprocity and palindromicity.
\end{proof}

A subset gives a partition inside a \(k\)-by-\((n-k)\) rectangle by taking
\(\lambda_j=i_{k+1-j}-(k-j)\). These integers decrease weakly, lie between
zero and \(n-k\), and sum to the exponent in \eqref{qcc:subset}. The construction
is reversible. This proves the rectangular-partition interpretation directly.
At \(q=1\), every subset has weight one, so \(\GB[1]nk=\binom nk\).

\begin{example}
\[
 \GB42=1+q+2q^2+q^3+q^4,\qquad
 \GB63=1+q+2q^2+3q^3+3q^4+3q^5+3q^6+2q^7+q^8+q^9.
\]
The coefficient sums are respectively \(6\) and \(20\); the symmetry centers
are respectively \(2\) and \(9/2\).
\end{example}

\section{The finite and reciprocal binomial theorems}
\begin{theorem}\label{qcc:finite-binomial}
For every \(N\ge0\),
\begin{equation}\label{qcc:finite}
 (z;q)_N=\sum_{k=0}^N(-1)^kq^{\binom{k}{2}}\GB Nk z^k.
\end{equation}
For \(N\ge1\), as a formal series in \(z\),
\begin{equation}\label{qcc:reciprocal}
 \frac1{(z;q)_N}=\sum_{r\ge0}\GB{N+r-1}r z^r.
\end{equation}
When \(N=0\), the reciprocal is one.
\end{theorem}
\begin{proof}
To obtain \(z^k\) in the finite product, choose the \(-zq^i\) term in precisely
\(k\) factors. Its weight is \((-1)^kq^{\sum i}\); comparison with
\eqref{qcc:subset} proves \eqref{qcc:finite}.

Write \(H_N(z)=1/(z;q)_N\). Then
\((1-q^{N-1}z)H_N(z)=H_{N-1}(z)\).
If \(c_{N,r}=[z^r]H_N\), it follows that
\(c_{N,r}=c_{N-1,r}+q^{N-1}c_{N,r-1}\).
The candidate \(\GB{N+r-1}r\) has this recurrence by
\eqref{qcc:pascal1}, with \(c_{1,r}=1\) and \(c_{N,0}=1\).
Induction on \(N+r\) proves the reciprocal identity.
\end{proof}

These classical identities are the finite and negative-power backbone of the
extension; compare the standard conventions in \cite{dlmf172,dlmf175}.

\section{A formal infinite-binomial identity without a topology error}
\begin{theorem}\label{qcc:infinite-binomial}
There is a unique \(R_a(z)\in\Q(q,a)[[z]]\) with \(R_a(0)=1\) satisfying
\begin{equation}\label{qcc:binomial-functional}
 (1-z)R_a(z)=(1-az)R_a(qz).
\end{equation}
It is
\begin{equation}\label{qcc:formal-binomial-series}
 R_a(z)=\sum_{r\ge0}\frac{(a;q)_r}{(q;q)_r}z^r.
\end{equation}
For complex \(|q|<1\) and \(|z|<1\), it also equals
\begin{equation}\label{qcc:analytic-binomial}
 R_a(z)=\frac{(az;q)_\infty}{(z;q)_\infty}.
\end{equation}
\end{theorem}
\begin{proof}
Comparison of coefficients in \eqref{qcc:binomial-functional} gives
\((1-q^r)c_r=(1-aq^{r-1})c_{r-1}\), with \(c_0=1\).
This proves both uniqueness and \eqref{qcc:formal-binomial-series}.
For \(|q|<1\), the infinite products converge locally uniformly in their
arguments, since the tails are bounded by convergent geometric sums. The
denominator has no zero for \(|z|<1\). Their ratio satisfies the same functional
equation, is analytic at zero, and has constant term one. Its Taylor coefficients
are therefore the coefficients already computed. Unless the series terminates, the coefficient ratio tends
to one, so the series converges absolutely for \(|z|<1\); termination
makes convergence immediate.
\end{proof}
\begin{auditbox}
With \(q\) merely an indeterminate, the product
\(\prod_{j\ge0}(1-zq^j)\) is not defined by \(z\)-adic convergence: its linear
coefficient changes at every step. Equation \eqref{qcc:formal-binomial-series}
or the functional equation is the formal definition. Equation
\eqref{qcc:analytic-binomial} is an additional analytic realization.
\end{auditbox}

\chapter{Convolution, inversion, and linearization}
\label{qcc:convolution}
\section{Two-block and many-block Vandermonde identities}
\begin{theorem}\label{qcc:vandermonde}
For \(m,n,r\ge0\),
\begin{equation}\label{qcc:vandermonde-formula}
 \GB{n+m}r=
 \sum_{k=\max(0,r-m)}^{\min(n,r)}
 q^{(n-k)(r-k)}\GB nk\GB m{r-k}.
\end{equation}
More generally, if \(n_1+\cdots+n_s=N\), then
\begin{equation}\label{qcc:block-vandermonde}
 \GB Nr=\sum_{\substack{0\le k_i\le n_i\\\sum k_i=r}}
 q^{\sum_{i<j}(n_i-k_i)k_j}\prod_{i=1}^s\GB{n_i}{k_i}.
\end{equation}
\end{theorem}
\begin{proof}
Split \((-z;q)_{n+m}=(-z;q)_n(-zq^n;q)_m\) and use
\eqref{qcc:finite} with \(-z\) in place of \(z\).
After dividing the coefficient of \(z^r\) by \(q^{\binom r2}\), the exponent is
\[
 \binom k2+\binom{r-k}2+n(r-k)-\binom r2=(n-k)(r-k).
\]
The same computation for consecutive blocks of lengths \(n_i\) gives
\(\sum_i\binom{k_i}2+\sum_{i<j}n_i k_j-\binom r2
 =\sum_{i<j}(n_i-k_i)k_j\), proving the general case.
\end{proof}
In particular, symmetry and the substitution \(k\mapsto n-k\) give
\begin{equation}\label{qcc:central}
 \sum_{k=0}^n q^{k^2}\GB nk^2=\GB{2n}n.
\end{equation}
The exponent \(k^2\) is essential; removing it does not give a Gaussian analogue
of the ordinary central-binomial identity.

\section{Gaussian inversion and the subspace lattice}
\begin{theorem}[Gaussian binomial inversion]\label{qcc:inversion}
For arbitrary sequences over \(\Q(q)\),
\begin{equation}\label{qcc:inverse-pair}
 b_n=\sum_{k=0}^n\GB nk a_k
 \quad\Longleftrightarrow\quad
 a_n=\sum_{k=0}^n(-1)^{n-k}q^{\binom{n-k}{2}}\GB nk b_k.
\end{equation}
Both identities also make sense over \(\Z[q]\) when the sequences do.
\end{theorem}
\begin{proof}
The coefficient of \(a_j\) after substitution on the right is
\[
 \sum_{k=j}^n(-1)^{n-k}q^{\binom{n-k}{2}}\GB nk\GB kj.
\]
Cancel factorials to obtain
\(\GB nk\GB kj=\GB nj\GB{n-j}{k-j}\).
The remaining sum is \((1;q)_{n-j}\), after reversing its index and using
symmetry. It is zero unless \(n=j\), when it is one. The two triangular
matrices therefore multiply to the identity. Their unit diagonal makes either
one-sided inverse also the inverse on every finite truncation.
\end{proof}

The same inverse has a geometric interpretation. For a prime power \(Q\), the
number of \(k\)-dimensional subspaces of \(\F_Q^n\) is \(\GB[Q]nk\).
Indeed, count ordered linearly independent \(k\)-tuples and divide by the number
of ordered bases of \(\F_Q^k\):
\[
 \frac{\prod_{j=0}^{k-1}(Q^n-Q^j)}{\prod_{j=0}^{k-1}(Q^k-Q^j)}=\GB[Q]nk.
\]
An interval of rank difference \(r\) in the subspace lattice is itself a
subspace lattice of dimension \(r\). Its incidence-algebra M\"obius function is
\((-1)^rQ^{\binom r2}\), because its defining convolution with the constant-one
function is exactly \((1;Q)_r=0\) for \(r>0\). This gives a direct counterpart
of the source article's use of incidence algebras for ordinary inversion.

\section{A positive linearization formula}
\begin{theorem}\label{qcc:linearization}
For \(0\le k,\ell\le n\),
\begin{equation}\label{qcc:linearization-formula}
 \GB nk\GB n\ell
 =\sum_{r=\max(k,\ell)}^{\min(n,k+\ell)}
 q^{(r-k)(r-\ell)}\GB nr\GB rk\GB k{k+\ell-r}.
\end{equation}
\end{theorem}
\begin{proof}
First let \(q=Q\) be a prime power and count ordered pairs \((V,W)\) of
subspaces of dimensions \(k,\ell\). Set \(U=V+W\), \(\dim U=r\), and
\(J=V\cap W\), \(\dim J=j=k+\ell-r\). Choose \(U\), then \(V\subseteq U\),
then \(J\subseteq V\), giving the three Gaussian factors.
In \(U/J\), the space \(V/J\) has dimension \(k-j=r-\ell\).
The possible \(W/J\) are complements of this space, of dimension
\(\ell-j=r-k\). Relative to one fixed complement, every other complement is
the graph of a unique linear map into \(V/J\). There are
\(Q^{(r-k)(r-\ell)}\) such maps. This proves the formula for every prime power.
Both sides are polynomials in \(q\); infinitely many equal values prove the
polynomial identity.
\end{proof}
For example,
\begin{equation}\label{qcc:qinteger-square}
 \qi n^2=\qi n+q(1+q)\GB n2.
\end{equation}
This is a Gaussian factorial-moment identity. A general power analogue follows
from \(q\)-Stirling basis change in \cref{qcc:qstirling}.

\section{Pochhammer convolution and terminating hypergeometric sums}
\begin{theorem}\label{qcc:poch-convolution}
For \(N\ge0\),
\begin{equation}\label{qcc:poch-convolution-formula}
 (ab;q)_N=\sum_{j=0}^N\GB Nj a^j(b;q)_j(a;q)_{N-j}.
\end{equation}
Consequently, wherever the displayed denominators are nonzero,
\begin{align}
 \sum_{j=0}^N
 \frac{(q^{-N};q)_j(a;q)_j}{(q;q)_j(c;q)_j}q^j
 &=a^N\frac{(c/a;q)_N}{(c;q)_N},\label{qcc:chu1}\\
 \sum_{j=0}^N
 \frac{(q^{-N};q)_j(a;q)_j}{(q;q)_j(c;q)_j}
 \left(\frac{cq^N}{a}\right)^j
 &=\frac{(c/a;q)_N}{(c;q)_N}.\label{qcc:chu2}
\end{align}
The last two formulae are the terminating \(q\)-Chu--Vandermonde identities.
\end{theorem}
\begin{proof}
The functional equation \eqref{qcc:binomial-functional} implies
\(R_{ab}(z)=R_a(z)R_b(az)\): both sides have constant term one and satisfy
\((1-z)R(z)=(1-abz)R(qz)\).
Comparison of \(z^N\) in \eqref{qcc:formal-binomial-series}, followed by
multiplication by \((q;q)_N\), proves \eqref{qcc:poch-convolution-formula}.
This proof is entirely formal.

Divide that formula by \((a;q)_N\). Direct reversal of finite products gives
\begin{equation}\label{qcc:conversion-chu}
 \GB Nj\frac{(a;q)_{N-j}}{(a;q)_N}a^j
 =\frac{(q^{-N};q)_j}{(q;q)_j(q^{1-N}/a;q)_j}q^j.
\end{equation}
For completeness, use
\[
 (q^{-N};q)_j=(-1)^j q^{-Nj+\binom j2}\frac{(q;q)_N}{(q;q)_{N-j}}
\]
and
\[
 (aq^{N-j};q)_j=(-a)^j q^{Nj-j(j+1)/2}(q^{1-N}/a;q)_j
\]
in the left side. The powers reduce to \(q^j\), proving
\eqref{qcc:conversion-chu}.
Now set the original \(a=q^{1-N}/c\) and rename the original \(b\) as \(a\).
The product reversal
\((x;q)_N=(-x)^Nq^{\binom N2}(q^{1-N}/x;q)_N\)
turns \((ab;q)_N/(a;q)_N\) into the right side of \eqref{qcc:chu1}.

Finally replace \(q,a,c\) in \eqref{qcc:chu1} by
\(q^{-1},a^{-1},c^{-1}\). Reversing each finite product changes the summand's
extra power into \((cq^N/a)^j\); the prefactor \(a^{-N}\) on the right cancels
the reversal factor \(a^N\). This yields \eqref{qcc:chu2}.
All calculations first take place in a rational-function field; any admissible
specialization then follows.
\end{proof}
The identities are classical; \cite{dlmf176} records their hypergeometric
normalizations. The proof above explains why they are coefficient-convolution
identities before they are hypergeometric summations.

\chapter{Symmetric functions and the Bell master formula}
\label{qcc:bell-master}
\section{Elementary, complete, and power-sum coordinates}
For a finite alphabet \(X=(x_1,\ldots,x_N)\), define
\[
 e_r(X)=\sum_{i_1<\cdots<i_r}x_{i_1}\cdots x_{i_r},\qquad
 h_r(X)=\sum_{i_1\le\cdots\le i_r}x_{i_1}\cdots x_{i_r},
 \qquad p_r(X)=\sum_i x_i^r.
\]
Both degree-zero functions are one. Direct multiplication gives
\begin{align}
 \sum_{r\ge0}e_r(X)z^r&=\prod_i(1+x_i z),\label{qcc:e-gf}\\
 \sum_{r\ge0}h_r(X)z^r&=\prod_i(1-x_i z)^{-1}.\label{qcc:h-gf}
\end{align}
Taking formal logarithms gives, for \(r\ge0\),
\begin{align}
 e_r(X)&=\frac1{r!}\Y r\bigl(p_1,-1!p_2,2!p_3,\ldots,
                         (-1)^{r-1}(r-1)!p_r\bigr),\label{qcc:e-bell}\\
 h_r(X)&=\frac1{r!}\Y r\bigl(p_1,1!p_2,2!p_3,\ldots,(r-1)!p_r\bigr).
 \label{qcc:h-bell}
\end{align}
These identities follow from \eqref{qcc:bell-def}; no additional combinatorial
normalization is involved. Also
\(\sum_{j=0}^r(-1)^j e_j(X)h_{r-j}(X)=\delta_{r0}\), because the two
generating functions with opposite signs multiply to one.

\begin{theorem}[Principal specialization]\label{qcc:principal}
For \(X_N=(1,q,\ldots,q^{N-1})\),
\begin{equation}\label{qcc:principal-formula}
 p_r(X_N)=\qi[q^r]N,\qquad
 e_k(X_N)=q^{\binom k2}\GB Nk,\qquad
 h_k(X_N)=\GB{N+k-1}k\quad(N\ge1).
\end{equation}
\end{theorem}
\begin{proof}
The first formula is a geometric sum. The other two follow by comparing
\eqref{qcc:e-gf} and \eqref{qcc:h-gf} with \eqref{qcc:finite} and
\eqref{qcc:reciprocal}.
\end{proof}
Thus a Gaussian coefficient is an ordinary symmetric polynomial at geometric
nodes. In particular,
\begin{align}
 q^{\binom k2}\GB Nk
 &=\frac1{k!}\Y k\left(\left((-1)^{r-1}(r-1)!\qi[q^r]N\right)_{r=1}^k\right),
 \label{qcc:gaussian-e-bell}\\
 \GB{N+k-1}k
 &=\frac1{k!}\Y k\left(\left((r-1)!\qi[q^r]N\right)_{r=1}^k\right).
 \label{qcc:gaussian-h-bell}
\end{align}
For instance,
\[
 \GB{N+1}2=\frac{\qi N^2+\qi[q^2]N}{2},\qquad
 q\GB N2=\frac{\qi N^2-\qi[q^2]N}{2}.
\]
The apparent halves cancel because these are symmetric-function identities over
\(\Z[q]\).

\section{Arbitrary powers of geometric products}
\begin{theorem}[Geometric-product coefficient formula]\label{qcc:master-theorem}
Let \(N_\ell\ge0\), \(d_\ell\ge1\), and let \(a_\ell,\alpha_\ell\) be
commuting parameters in a \(\Q\)-algebra. Define formally
\begin{equation}\label{qcc:master-product}
 F(z)=\prod_{\ell=1}^L(a_\ell z;q^{d_\ell})_{N_\ell}^{-\alpha_\ell},
 \qquad
 P_r=\sum_{\ell=1}^L\alpha_\ell a_\ell^r\qi[q^{d_\ell r}]{N_\ell}.
\end{equation}
Powers with nonintegral exponents mean \(\exp(-\alpha_\ell\log(\cdot))\)
with constant term one. Then
\begin{equation}\label{qcc:master-coefficient}
 [z^n]F(z)=\frac1{n!}\Y n\bigl(P_1,1!P_2,\ldots,(n-1)!P_n\bigr)
 =\sum_{\substack{m_1,\ldots,m_n\ge0\\\sum r m_r=n}}
 \prod_{r=1}^n\frac{P_r^{m_r}}{r^{m_r}m_r!}.
\end{equation}
\end{theorem}
\begin{proof}
Each factor has a formal logarithm, and
\[
 \log F(z)
 =\sum_{\ell=1}^L\alpha_\ell\sum_{j=0}^{N_\ell-1}
       \sum_{r\ge1}\frac{(a_\ell q^{d_\ell j}z)^r}{r}
 =\sum_{r\ge1}P_r\frac{z^r}{r}.
\]
At every fixed degree this rearranges only finitely many contributions from the
finite alphabets. Apply \eqref{qcc:bell-def} with
\(x_r=(r-1)!P_r\), then \eqref{qcc:bell-finite}.
\end{proof}
The first four coefficients are
\begin{align*}
 f_1&=P_1,\qquad f_2=\frac{P_1^2+P_2}{2},\\
 f_3&=\frac{P_1^3+3P_1P_2+2P_3}{6},\\
 f_4&=\frac{P_1^4+6P_1^2P_2+3P_2^2+8P_1P_3+6P_4}{24}.
\end{align*}
Negative integral \(\alpha_\ell\) encode numerator products, positive integral
ones encode repeated reciprocal products, and arbitrary parameters encode
binomial powers. Several bases \(q^{d_\ell}\) require no new machinery.

\begin{corollary}[A two-base convolution]\label{qcc:two-base}
For \(M,N\ge1\),
\begin{multline}\label{qcc:two-base-formula}
 \sum_{j=0}^n a^j b^{n-j}
 \GB[q]{M+j-1}j\GB[q^d]{N+n-j-1}{n-j}\\
 =\frac1{n!}\Y n\left(
 \left((r-1)!\bigl(a^r\qi[q^r]M+b^r\qi[q^{dr}]N\bigr)\right)_{r=1}^n
 \right).
\end{multline}
\end{corollary}
\begin{proof}
Expand the two reciprocal products using \eqref{qcc:reciprocal} and multiply.
Alternatively, apply \cref{qcc:master-theorem} directly to their product. Equating
the two coefficients proves the identity.
\end{proof}

\section{A determinant and the recovery of logarithmic data}
\begin{proposition}\label{qcc:determinant}
If \(F(z)=\exp(\sum_{r\ge1}P_rz^r/r)=\sum f_nz^n\), then
\begin{equation}\label{qcc:det-formula}
 f_n=\frac1{n!}
 \det\begin{pmatrix}
 P_1&-1&0&\cdots&0\\
 P_2&P_1&-2&\cdots&0\\
 P_3&P_2&P_1&\ddots&\vdots\\
 \vdots&\vdots&\ddots&\ddots&-(n-1)\\
 P_n&P_{n-1}&\cdots&P_2&P_1
 \end{pmatrix}.
\end{equation}
Conversely, put \(x_j=j!f_j\). Then
\begin{equation}\label{qcc:inverse-bell-log}
 P_n=\frac1{(n-1)!}\sum_{k=1}^n(-1)^{k-1}(k-1)!
                   \B nk(x_1,\ldots,x_{n-k+1}).
\end{equation}
\end{proposition}
\begin{proof}
The equation \(F'=(\sum_{r\ge1}P_rz^{r-1})F\) gives
\(nf_n=\sum_{r=1}^nP_rf_{n-r}\). If \(D_n\) is the determinant above, expansion
along its last row gives
\[
 D_n=\sum_{r=1}^n\frac{(n-1)!}{(n-r)!}P_rD_{n-r},\qquad D_0=1.
\]
The superdiagonal products supply the factorial ratio, and the Laplace signs
cancel their signs. Therefore \(D_n/n!\) satisfies the same recurrence and
initial value as \(f_n\).
For the inverse formula, expand
\(\log F=\sum_{k\ge1}(-1)^{k-1}(F-1)^k/k\), use
\eqref{qcc:partial-def}, and compare \(z^n\) with \(P_n/n\).
\end{proof}
This is precisely the moment--cumulant mechanism in the source article, with
geometric power sums serving as the logarithmic data.

\chapter{Partial fractions and divisor identities}
\label{qcc:divisor}
\section{An interpolation formula for complete symmetric functions}
\begin{lemma}\label{qcc:partial-fraction}
If \(x_1,\ldots,x_N\) are distinct and nonzero, then
\begin{equation}\label{qcc:partial-fraction-formula}
 \prod_{i=1}^N\frac1{1-x_i z}
 =\sum_{i=1}^N\frac{x_i^{N-1}}{\prod_{j\ne i}(x_i-x_j)}\frac1{1-x_i z}.
\end{equation}
Consequently,
\begin{equation}\label{qcc:h-divided}
 h_r(x_1,\ldots,x_N)
 =\sum_{i=1}^N\frac{x_i^{r+N-1}}{\prod_{j\ne i}(x_i-x_j)}.
\end{equation}
\end{lemma}
\begin{proof}
Multiply the proposed rational identity by \(\prod_j(1-x_jz)\).
Both sides are polynomials of degree at most \(N-1\); at \(z=1/x_i\), only
one summand on the right survives, and it equals one. Agreement at \(N\)
distinct points proves the identity. Expanding each geometric denominator and
comparing coefficients proves \eqref{qcc:h-divided}.
\end{proof}
The divided-difference expression on the right has removable singularities when
nodes coincide; the left side supplies the polynomial continuation. This is a
simple instance of why a rational summation formula can encode a polynomial.

\section{A finite divisor-sum master identity}
\begin{theorem}[Dilcher-type identity and Bell form]\label{qcc:dilcher}
For \(N\ge1\), \(r\ge0\), and indeterminate \(q\), let
\(u_j=q^j/(1-q^j)\). Then
\begin{align}
 &\sum_{j=1}^N(-1)^{j-1}
  \frac{q^{\binom j2+rj}}{(1-q^j)^r}\GB Nj
 =h_r(u_1,\ldots,u_N)\label{qcc:dilcher-formula}\\
 &\hspace{13mm}
 =\sum_{1\le i_1\le\cdots\le i_r\le N}
      \frac{q^{i_1+\cdots+i_r}}{(1-q^{i_1})\cdots(1-q^{i_r})}
 \nonumber\\
 &\hspace{13mm}
 =\frac1{r!}\Y r\left(
       \left((s-1)!\sum_{j=1}^N\frac{q^{js}}{(1-q^j)^s}\right)_{s=1}^r
       \right).
 \label{qcc:dilcher-bell}
\end{align}
For \(r=0\), every right-hand side is interpreted as one.
\end{theorem}
\begin{proof}
Apply \cref{qcc:partial-fraction} to \(q,q^2,\ldots,q^N\). For the node
\(q^j\), the residue coefficient is
\[
 \frac1{\prod_{i\ne j}(1-q^{i-j})}
 =\frac{(-1)^{j-1}q^{\binom j2}}{(q;q)_{j-1}(q;q)_{N-j}}.
\]
Multiplication by \((q;q)_N\) therefore yields
\begin{equation}\label{qcc:divisor-rational}
 \frac{(q;q)_N}{(zq;q)_N}
 =\sum_{j=1}^N(-1)^{j-1}q^{\binom j2}(1-q^j)\GB Nj\frac1{1-zq^j}.
\end{equation}
Set \(z=1+t\). The left side becomes
\(\prod_{j=1}^N(1-u_jt)^{-1}\).
The coefficient of \(t^r\) on the right is exactly the left side of
\eqref{qcc:dilcher-formula}. The complete-symmetric and Bell descriptions now
follow from \eqref{qcc:h-gf} and \eqref{qcc:h-bell}.
\end{proof}
This is the finite identity associated with Dilcher's divisor-function
identities \cite{dilcher}. The partial-fraction proof simultaneously explains the
alternating single sum, the weakly ordered multiple sum, and the Bell polynomial.
For \(r=1\), it reduces to
\[
 \sum_{j=1}^N(-1)^{j-1}\frac{q^{j(j+1)/2}}{1-q^j}\GB Nj
 =\sum_{j=1}^N\frac{q^j}{1-q^j}.
\]
For \(r=2\), the right side is \(\frac12((\sum u_j)^2+\sum u_j^2)\),
and for \(r=3\), it is
\(\frac16((\sum u_j)^3+3(\sum u_j)(\sum u_j^2)+2\sum u_j^3)\).

\section{The ordinary harmonic limit}
Multiply \eqref{qcc:dilcher-formula} by \((1-q)^r\) and let \(q\to1\).
Since \((1-q)u_j\to1/j\), the limit is
\begin{equation}\label{qcc:harmonic-limit}
 \sum_{j=1}^N(-1)^{j-1}\binom Nj\frac1{j^r}
 =h_r\left(1,\frac12,\ldots,\frac1N\right)
 =\frac1{r!}\Y r\bigl(H_N^{(1)},1!H_N^{(2)},\ldots,(r-1)!H_N^{(r)}\bigr),
\end{equation}
where \(H_N^{(s)}=\sum_{j=1}^N j^{-s}\). This limit is termwise because all sums
are finite. Thus the ordinary binomial--harmonic identities of coefficient
calculus are literal degenerations of the Gaussian identities, not merely
analogous-looking formulae.

\chapter{Jackson differentiation and geometric Newton interpolation}
\label{qcc:jackson}
\section{The operator and its commutation rule}
Let \(T_qf(x)=f(qx)\), and let \(X\) denote multiplication by \(x\). Define
\begin{equation}\label{qcc:jackson-def}
 D_qf(x)=\frac{f(x)-f(qx)}{(1-q)x}.
\end{equation}
For polynomials, this has no singularity at \(x=0\), because the numerator has
zero constant term. It acts by
\(D_qx^m=[m]_qx^{m-1}\), and hence
\begin{equation}\label{qcc:q-weyl}
 D_qX=1+qXD_q,\qquad D_qT_q=qT_qD_q.
\end{equation}
These identities follow on every monomial and hence on every polynomial. This is
the convention for the Jackson derivative used throughout; compare
\cite{dlmf172}.

\begin{theorem}[The shifted product rule]\label{qcc:leibniz}
For \(n\ge0\),
\begin{equation}\label{qcc:q-leibniz}
 D_q^n(fg)(x)=\sum_{k=0}^n\GB nk
 (D_q^k f)(q^{n-k}x)(D_q^{n-k}g)(x).
\end{equation}
\end{theorem}
\begin{proof}
The first-order formula is
\(D_q(fg)=(D_qf)g+(T_qf)D_qg\), obtained by adding and subtracting
\(f(qx)g(x)\) in the numerator. Suppose the formula is true at order \(n\).
Differentiating a summand gives a contribution with \(k\) unchanged and a
contribution with \(k\) increased by one. In the latter, differentiation of
\((D_q^kf)(q^{n-k}x)\) produces a factor \(q^{n-k}\).
Thus the new coefficient at index \(k\) is
\(\GB nk+q^{n-k+1}\GB n{k-1}=\GB{n+1}k\), by
\eqref{qcc:pascal1}. The shifts are exactly those in the stated formula.
\end{proof}
\begin{auditbox}
The argument shift \(q^{n-k}x\) is part of the theorem. Replacing the ordinary
binomial coefficient by a Gaussian coefficient in the usual Leibniz rule, while
leaving all factors evaluated at \(x\), is false.
\end{auditbox}

\section{A geometric Taylor basis}
Define
\begin{equation}\label{qcc:q-newton-basis}
 (x-a)_{q}^{[k]}=\prod_{j=0}^{k-1}(x-aq^j).
\end{equation}
This is a polynomial basis in \(x\), not a power of \(x-a\).
\begin{theorem}[Jackson--Newton expansion]\label{qcc:q-taylor}
For every polynomial \(p\) of degree at most \(d\),
\begin{equation}\label{qcc:q-taylor-formula}
 p(x)=\sum_{k=0}^d\frac{(D_q^kp)(a)}{\qf k}(x-a)_q^{[k]}.
\end{equation}
The identity first holds over \(\Q(q,a)\), and specializes wherever its
coefficients are defined. At \(q=1\) it becomes ordinary Taylor expansion at
\(a\).
\end{theorem}
\begin{proof}
The polynomial \(D_q(x-a)_q^{[k]}\) has degree \(k-1\) and leading coefficient
\([k]_q\). At each \(x=aq^j\), \(0\le j\le k-2\), its two numerator terms
vanish. Therefore
\[
 D_q(x-a)_q^{[k]}=[k]_q(x-a)_q^{[k-1]}.
\]
This argument at generic \(a,q\) proves a polynomial identity and hence also
covers \(a=0\). Expand \(p\) in this monic basis, apply \(D_q^k\), and set
\(x=a\). Terms with index less than \(k\) vanish under differentiation; terms
with larger index retain a factor \(x-a\); the index \(k\) term yields
\([k]_q!\) times its coefficient. This proves the formula.
\end{proof}
For example,
\[
 x^2=a^2+(1+q)a(x-a)+(x-a)(x-aq),
\]
and
\[
 x^3=a^3+[3]_qa^2(x-a)+[3]_qa(x-a)(x-aq)
                  +(x-a)(x-aq)(x-aq^2).
\]
At a sample node \(x=aq^n\), factors with \(k>n\) vanish, and direct
factorization gives the discrete Gaussian transform
\begin{equation}\label{qcc:sample-transform}
 p(aq^n)=\sum_{k=0}^n\GB nk\bigl(a(q-1)\bigr)^k
                  q^{\binom k2}(D_q^kp)(a).
\end{equation}
Indeed,
\(\prod_{j<k}(aq^n-aq^j)
 =[a(q-1)]^kq^{\binom k2}[k]_q!\GB nk\).
This identifies Gaussian inversion with the extraction of geometric Newton
coefficients.

\section{Explicit stencils and divided differences}
\begin{theorem}\label{qcc:stencil}
For \(n\ge0\),
\begin{equation}\label{qcc:q-stencil}
 D_q^nf(x)=\frac1{(1-q)^nx^n}
 \sum_{j=0}^n(-1)^j q^{\binom j2-(n-1)j}\GB nj f(q^jx).
\end{equation}
At distinct nonzero nodes \(a,aq,\ldots,aq^n\),
\begin{equation}\label{qcc:divided-difference}
 f[a,aq,\ldots,aq^n]=\frac{(D_q^nf)(a)}{\qf n}.
\end{equation}
The left side denotes the leading coefficient of the degree-at-most-\(n\)
interpolating polynomial through the given data.
\end{theorem}
\begin{proof}
The commutation relation \(T_qX^{-1}=q^{-1}X^{-1}T_q\) gives
\[
 D_q^n=(1-q)^{-n}X^{-n}\prod_{r=0}^{n-1}(1-q^{-r}T_q).
\]
Expand the finite product with base \(q^{-1}\), and use reciprocity
\eqref{qcc:reciprocity}. Its \(T_q^j\) coefficient becomes
\((-1)^jq^{\binom j2-(n-1)j}\GB nj\), proving
\eqref{qcc:q-stencil}.
Let \(p\) interpolate \(f\) at these nodes. Formula \eqref{qcc:q-stencil}
shows that \((D_q^nf)(a)=(D_q^np)(a)\). The latter equals \([n]_q!\) times
the leading coefficient of \(p\), either by monomial differentiation or by
\cref{qcc:q-taylor}. This proves \eqref{qcc:divided-difference}.
\end{proof}

\begin{corollary}[An exact Gaussian coefficient extractor]\label{qcc:extractor}
If \(p\) has degree at most \(n\), then
\begin{equation}\label{qcc:extractor-formula}
 [x^n]p(x)=\frac1{(q;q)_n}
 \sum_{j=0}^n(-1)^j q^{\binom j2-(n-1)j}\GB nj p(q^j).
\end{equation}
The numerator annihilates \(x^r\) for every \(0\le r<n\).
\end{corollary}
\begin{proof}
Apply \eqref{qcc:q-stencil} at \(x=1\) and divide by \([n]_q!\).
Alternatively, the numerator on \(x^r\) is
\((q^{r-n+1};q)_n\), by \eqref{qcc:finite}. One factor is zero for
\(0\le r<n\), while for \(r=n\) the product is \((q;q)_n\).
\end{proof}
This is the geometric counterpart of the additive finite-difference stencil in
the source. Its cancelling objects are monomials evaluated at geometric nodes.

\chapter{Two \qbname-Stirling arrays and normal ordering}
\label{qcc:qstirling}
\section{The monic \qbname-integer factorial basis}
Define
\begin{equation}\label{qcc:qstirling-basis}
 P_k(x)=\prod_{j=0}^{k-1}(x-[j]_q),\qquad
 P_n(x)=\sum_{k=0}^n s_q(n,k)x^k,\qquad
 x^n=\sum_{k=0}^n S_q(n,k)P_k(x).
\end{equation}
These are changes between two monic polynomial bases. They define \(s_q,S_q\)
unambiguously over \(\Z[q]\), even when specialized nodes coincide.

\begin{theorem}\label{qcc:qstirling-basic}
The arrays are mutually inverse, and
\begin{align}
 s_q(n+1,k)&=s_q(n,k-1)-[n]_qs_q(n,k),\label{qcc:sq-rec}\\
 S_q(n+1,k)&=S_q(n,k-1)+[k]_qS_q(n,k),\label{qcc:Sq-rec}\\
 s_q(n,k)&=(-1)^{n-k}e_{n-k}([0]_q,\ldots,[n-1]_q),\label{qcc:sq-e}\\
 S_q(n,k)&=h_{n-k}([0]_q,\ldots,[k]_q).\label{qcc:Sq-h}
\end{align}
Both arrays have entry one at \((0,0)\) and vanish outside the triangular range.
\end{theorem}
\begin{proof}
Multiplication of \(P_n(x)\) by \(x-[n]_q\) gives
\eqref{qcc:sq-rec}. The identity \(xP_k=P_{k+1}+[k]_qP_k\) gives
\eqref{qcc:Sq-rec}. The coefficient formula \eqref{qcc:sq-e} follows by
choosing roots from the product defining \(P_n\).
For \eqref{qcc:Sq-h}, use
\[
 h_r([0]_q,\ldots,[k]_q)
 =h_r([0]_q,\ldots,[k-1]_q)+[k]_qh_{r-1}([0]_q,\ldots,[k]_q).
\]
With \(r=n-k\), this is precisely the recurrence for \(S_q\), with the same
boundary values. Finally, substituting either basis change into the other and
using uniqueness of coefficients proves matrix inversion.
\end{proof}
At \(q=1\), these are the ordinary signed first-kind and second-kind Stirling
numbers. At \(q=0\), all positive-index nodes equal one, so for \(n\ge1\),
\[
 S_0(n,k)=\binom{n-1}{k-1},\qquad
 s_0(n,k)=(-1)^{n-k}\binom{n-1}{k-1}\quad(1\le k\le n).
\]
This follows from \(P_n(x)=x(x-1)^{n-1}\) and from the complete-symmetric
formula. It illustrates the usefulness of a polynomial definition at degenerate
bases.

\section{A Gaussian closed sum, with every power of \qbname accounted for}
\begin{theorem}\label{qcc:qstirling-explicit}
Over \(\Q(q)\), for \(0\le k\le n\),
\begin{equation}\label{qcc:qstirling-sum}
 S_q(n,k)=\frac{q^{-\binom k2}}{[k]_q!}
 \sum_{j=0}^k(-1)^{k-j}q^{\binom{k-j}{2}}\GB kj[j]_q^n.
\end{equation}
The expression has the polynomial continuation specified by
\eqref{qcc:qstirling-basis}.
\end{theorem}
\begin{proof}
The coefficient of the degree-\(k\) Newton basis polynomial in the interpolation
of \(x^n\) at \([0]_q,\ldots,[k]_q\) is
\[
 \sum_{j=0}^k\frac{[j]_q^n}{\prod_{0\le i\le k,\ i\ne j}([j]_q-[i]_q)}.
\]
It equals the coefficient in the full Newton expansion because terms of degree
larger than \(k\) vanish at all these nodes. For \(i<j\),
\([j]_q-[i]_q=q^i[j-i]_q\), whereas for \(i>j\) it is
\(-q^j[i-j]_q\). Hence the denominator is
\[
 (-1)^{k-j}q^{jk-j(j+1)/2}[j]_q![k-j]_q!.
\]
Its inverse, multiplied by \([k]_q!\), has exponent
\(\binom{j+1}{2}-jk=\binom{k-j}{2}-\binom k2\).
This proves \eqref{qcc:qstirling-sum}.
\end{proof}
There is no guesswork in the factor \(q^{-\binom k2}\): it is the normalization
of the monic \(q\)-integer Newton basis. Omitting it produces a different array.

\section{Bell formulae for every diagonal}
For \(k\ge0\) and \(r\ge1\), let
\begin{equation}\label{qcc:qinteger-powersum}
 U_r(k)=\sum_{j=0}^k[j]_q^r
 =\frac1{(1-q)^r}\sum_{a=0}^r(-1)^a\binom ra[k+1]_{q^a}.
\end{equation}
The \(a=0\) term means \(k+1\). The equality follows by expanding
\((1-q^j)^r\) binomially and summing the geometric powers \(q^{aj}\).
Equations \eqref{qcc:sq-e}--\eqref{qcc:Sq-h} and the Bell formulae give
\begin{align}
 S_q(n,k)&=\frac1{(n-k)!}\Y{n-k}
   \left(\left((r-1)!U_r(k)\right)_{r=1}^{n-k}\right),\label{qcc:Sq-bell}\\
 s_q(n,k)&=\frac{(-1)^{n-k}}{(n-k)!}\Y{n-k}
   \left(\left((-1)^{r-1}(r-1)!U_r(n-1)\right)_{r=1}^{n-k}\right).
 \label{qcc:sq-bell}
\end{align}
These are recurrence-free finite formulae involving only binomial coefficients,
\(q\)-integers, factorials, and the finite sum \eqref{qcc:bell-finite}.
For example,
\[
 S_q(n,n-1)=\sum_{j=1}^{n-1}[j]_q
          =\sum_{a=0}^{n-2}(n-1-a)q^a,
\]
and
\[
 S_q(n,n-2)=\frac12\left(U_1(n-2)^2+U_2(n-2)\right).
\]
The latter uses the alphabet ending at \([n-2]_q\), not at \([n-1]_q\).

\section{The normal-ordering normalization}
Define
\begin{equation}\label{qcc:normalized-stirling}
 \widetilde S_q(n,k)=q^{\binom k2}S_q(n,k).
\end{equation}
Then
\begin{equation}\label{qcc:normalized-stirling-rec}
 \widetilde S_q(n+1,k)=[k]_q\widetilde S_q(n,k)
                         +q^{k-1}\widetilde S_q(n,k-1).
\end{equation}
Both this and \eqref{qcc:Sq-rec} occur in \(q\)-Stirling literature; the
restricted-growth-word viewpoint and related identities are discussed in
\cite{cai}. The defining recurrence must accompany any comparison of conventions.

\begin{theorem}[Normal ordering and factorial moments]\label{qcc:normal-ordering}
On polynomials,
\begin{equation}\label{qcc:normal-ordering-formula}
 (XD_q)^n=\sum_{k=0}^n\widetilde S_q(n,k)X^kD_q^k.
\end{equation}
For integers \(m,n\ge0\),
\begin{equation}\label{qcc:factorial-moments}
 [m]_q^n=\sum_{k=0}^{\min(m,n)}
          \widetilde S_q(n,k)[k]_q!\GB mk.
\end{equation}
\end{theorem}
\begin{proof}
Induction from \eqref{qcc:q-weyl} gives
\(D_qX^k=[k]_qX^{k-1}+q^kX^kD_q\).
Left multiplication of the right side of \eqref{qcc:normal-ordering-formula}
by \(XD_q\) therefore produces recurrence
\eqref{qcc:normalized-stirling-rec}, with the correct initial value.
This proves normal ordering.
Apply both sides to \(x^m\). The left side is \([m]_q^nx^m\), and
\(X^kD_q^kx^m=[k]_q!\GB mk x^m\); terms with \(k>m\) vanish.
Comparison proves \eqref{qcc:factorial-moments}.
\end{proof}
For \(n=2\), normal ordering reads
\((XD_q)^2=XD_q+qX^2D_q^2\). Thus
\(\widetilde S_q(2,2)=q\), while \(S_q(2,2)=1\). This tiny example is an
effective audit of the two normalizations.

\chapter{Weighted Bell polynomials and Riordan composition}
\label{qcc:composition}
\section{Gaussian convolution as a change of coefficient coordinates}
Set \(w_n=[n]_q!\), and record a sequence by
\[
 A(z)=\sum_{n\ge0}a_n\frac{z^n}{w_n}.
\]
Ordinary multiplication of series is then the Gaussian convolution
\begin{equation}\label{qcc:q-convolution}
 (a\star_q b)_n=\sum_{k=0}^n\GB nk a_kb_{n-k}.
\end{equation}
Associativity and commutativity follow from ordinary series multiplication.
More explicitly, for \(s\) factors the convolution coefficient is
\begin{equation}\label{qcc:q-multinomial}
 \sum_{i_1+\cdots+i_s=n}
 \frac{[n]_q!}{[i_1]_q!\cdots[i_s]_q!}
 a^{(1)}_{i_1}\cdots a^{(s)}_{i_s}.
\end{equation}
The factorial ratio is the Gaussian multinomial. It is polynomial because it is
the product
\(\GB n{i_1}\GB{n-i_1}{i_2}\cdots\).
Thus associativity extends polynomially to every specialization of \(q\), even
when the divided-power representation itself is not available.

\begin{proposition}[Explicit convolution inverse]\label{qcc:convolution-inverse}
If \(a_0\) is invertible, the inverse \(b\) for \(\star_q\) satisfies
\(b_0=a_0^{-1}\), and for \(n>0\),
\begin{equation}\label{qcc:convolution-inverse-formula}
 b_n=\frac1{a_0}\sum_{r=1}^n\frac{(-1)^r}{a_0^r}
 \sum_{\substack{i_1+\cdots+i_r=n\\i_j\ge1}}
 \frac{[n]_q!}{[i_1]_q!\cdots[i_r]_q!}\prod_{j=1}^r a_{i_j}.
\end{equation}
\end{proposition}
\begin{proof}
Write \(A(z)=a_0(1+U(z))\), where \(U(0)=0\). The formal geometric expansion
\(A(z)^{-1}=a_0^{-1}\sum_{r\ge0}(-U(z))^r\) is coefficientwise finite.
Multiplying its degree-\(n\) coefficient by \(w_n\) gives the displayed formula.
\end{proof}
In particular, with
\[
 e_q(z)=\sum_{n\ge0}\frac{z^n}{[n]_q!},\qquad
 E_q(z)=\sum_{n\ge0}q^{\binom n2}\frac{z^n}{[n]_q!},
\]
Gaussian inversion is the coefficient identity \(e_q(z)E_q(-z)=1\).

\section{Composition coefficients without an ambiguous ``\qbname-Bell'' convention}
Let
\(B(z)=\sum_{j\ge1}b_jz^j/w_j\). Define
\begin{equation}\label{qcc:weighted-bell-def}
 \qB nk(b)=\frac{w_n}{w_k}[z^n]B(z)^k.
\end{equation}
This is a specified weighted coefficient family; it is not a claim that all
objects called \(q\)-Bell polynomials in the literature have this meaning.

\begin{theorem}\label{qcc:weighted-bell}
For \(0\le k\le n\),
\begin{equation}\label{qcc:weighted-bell-finite}
 \qB nk(b)=\frac{w_nk!}{w_k}
 \sum_{\substack{m_1,\ldots,m_n\ge0\\\sum m_j=k,\ \sum j m_j=n}}
 \prod_{j=1}^n\frac1{m_j!}\left(\frac{b_j}{w_j}\right)^{m_j}.
\end{equation}
Equivalently,
\begin{equation}\label{qcc:weighted-bell-transport}
 \qB nk(b)=\frac{w_nk!}{w_kn!}\B nk
       \left(\frac{1!b_1}{w_1},\frac{2!b_2}{w_2},\ldots\right).
\end{equation}
If \(A(z)=\sum_{k\ge0}a_kz^k/w_k\), then
\begin{equation}\label{qcc:composition-formula}
 w_n[z^n]A(B(z))=\sum_{k=0}^n a_k\qB nk(b).
\end{equation}
\end{theorem}
\begin{proof}
Expand \(B(z)^k\) by the ordinary multinomial theorem. A multiplicity vector
\((m_j)\) contributes \(k!\prod_j(b_j/w_j)^{m_j}/m_j!\).
The two constraints count the factors and their total degree. This proves
\eqref{qcc:weighted-bell-finite}. Compare it with the defining finite sum for
\(\B nk\) to obtain \eqref{qcc:weighted-bell-transport}. Finally, expand
\(A(B)\) in powers of \(B\) and use \eqref{qcc:weighted-bell-def}.
\end{proof}
At \(q=1\), \(w_n=n!\), and \(\mathfrak B^{[q]}_{n,k}\) reduces to the
ordinary exponential Bell polynomial \(\B nk\). At \(q=0\), \(w_n=1\),
and it becomes the ordinary coefficient \([z^n]B(z)^k\). This cleanly separates
coefficient normalization from the choice of a deformation of set partitions.

\section{Weighted Riordan arrays}
For \(A(0)\ne0\) and \(B(0)=0\) with invertible linear coefficient, put
\begin{equation}\label{qcc:riordan-def}
 \mathcal M_q(A,B)_{n,k}=\frac{w_n}{w_k}[z^n]A(z)B(z)^k.
\end{equation}
\begin{theorem}\label{qcc:riordan-group}
These lower triangular matrices obey
\begin{align}
 \mathcal M_q(A,B)\mathcal M_q(C,D)
 &=\mathcal M_q\bigl(A\,C(B),D(B)\bigr),\label{qcc:riordan-product}\\
 \mathcal M_q(A,B)^{-1}
 &=\mathcal M_q\left(\frac1{A(\overline B)},\overline B\right),
 \label{qcc:riordan-inverse}
\end{align}
where \(\overline B\) is the ordinary compositional inverse of \(B\).
\end{theorem}
\begin{proof}
The \((n,k)\) entry of the product equals
\[
 \frac{w_n}{w_k}\sum_j[z^n]A(z)B(z)^j\,[t^j]C(t)D(t)^k.
\]
For fixed \(n\), the sum is finite because \(B(0)=0\). Summing before taking
\([z^n]\) produces \(A(z)C(B(z))D(B(z))^k\), proving the product law.
Inserting \(D=\overline B\) and \(C=1/A(\overline B)\) gives
\(\mathcal M_q(1,z)\), the identity matrix. The same substitution in the other
order gives the other inverse identity.
\end{proof}
This is diagonal transport of ordinary Riordan calculus. It extends the source's
Riordan framework while making clear which operation is unchanged: composition
is still ordinary substitution of formal power series.

\section{Why Jackson composition is a different problem}
The first-order chain rule for the Jackson derivative is
\begin{equation}\label{qcc:divided-chain}
 D_q(F\circ G)(x)=F[G(x),G(qx)]\,D_qG(x),
\end{equation}
where \(F[u,v]=(F(u)-F(v))/(u-v)\), with its polynomial continuation on the
diagonal. This follows simply by factoring the numerator in
\eqref{qcc:jackson-def}. It is not usually \((D_qF)(G(x))D_qG(x)\).
For example, with \(F(u)=u^2\) and \(G(x)=x^2\), the actual answer is
\([4]_qx^3\), whereas that guess gives \([2]_q^2x^3\).
The weighted Bell family in this chapter describes ordinary composition and
must not be silently substituted into a higher Jackson chain rule.

\chapter{Lagrange inversion with geometric product data}
\label{qcc:lagrange}
\section{The ordinary inversion theorem over a \qbname-dependent field}
\begin{theorem}[Formal Lagrange--B\"urmann]\label{qcc:lagrange-theorem}
Let \(\Phi(0)=1\). There is a unique \(y(z)\in zK[[z]]\) satisfying
\(y=z\Phi(y)\). For \(n\ge1\),
\begin{equation}\label{qcc:lagrange-formula}
 [z^n]H(y(z))=\frac1n[t^{n-1}]H'(t)\Phi(t)^n.
\end{equation}
In particular, for \(s\ge1\) and \(n\ge s\),
\begin{equation}\label{qcc:lagrange-powers}
 [z^n]y(z)^s=\frac{s}{n}[t^{n-s}]\Phi(t)^n.
\end{equation}
\end{theorem}
\begin{proof}
The equation determines the coefficient of degree \(n\) of \(y\) from lower
degrees, starting with \([z]y=1\). This proves existence and uniqueness.
The inverse substitution is \(z=t/\Phi(t)\).
For a Laurent series, residue means the coefficient of degree \(-1\).
The residue of a formal derivative is zero, and the substitution rule
\(\Res_z A(u(z))u'(z)\,dz=\Res_u A(u)\,du\) follows on monomials:
for powers other than \(-1\), the integrand is a derivative; for the power
\(-1\), the logarithmic derivative has residue one.
Thus
\begin{align*}
 [z^n]H(y(z))
 &=\Res_t H(t)\left(\frac{\Phi(t)^n}{t^{n+1}}
                   -\frac{\Phi(t)^{n-1}\Phi'(t)}{t^n}\right)\dd t\\
 &=\frac1n\Res_t H'(t)\frac{\Phi(t)^n}{t^n}\dd t.
\end{align*}
The second equality is obtained by taking the zero residue of the derivative
of \(H(t)\Phi(t)^nt^{-n}\). It is \eqref{qcc:lagrange-formula}; taking
\(H(t)=t^s\) proves \eqref{qcc:lagrange-powers}.
\end{proof}
This is the ordinary theorem used in the source, now over a field containing
\(q\); see also the proof-oriented survey \cite{gessel}.

\section{A recurrence-free inverse for a product of shifted factorials}
\begin{theorem}\label{qcc:inverse-product}
Suppose \(\Phi(t)\) is the product \eqref{qcc:master-product}, with logarithmic
power sums \(P_r\). The solution of \(y=z\Phi(y)\) has coefficients
\begin{equation}\label{qcc:inverse-product-formula}
 [z^n]y(z)=\frac1{n!}\Y{n-1}
       \left(\left(n(r-1)!P_r\right)_{r=1}^{n-1}\right)\qquad(n\ge1).
\end{equation}
More generally,
\begin{equation}\label{qcc:inverse-product-powers}
 [z^n]y(z)^s=\frac{s}{n(n-s)!}\Y{n-s}
       \left(\left(n(r-1)!P_r\right)_{r=1}^{n-s}\right),\qquad n\ge s\ge1.
\end{equation}
\end{theorem}
\begin{proof}
The logarithm of \(\Phi(t)^n\) has power sums \(nP_r\).
Apply \cref{qcc:master-theorem} inside \eqref{qcc:lagrange-powers}.
The case \(s=1\) has denominator \(n(n-1)!=n!\).
\end{proof}
This solves, coefficient by coefficient and without auxiliary recurrence, the
inverse problem for
\[
 f(x)=x\prod_{\ell=1}^L(a_\ell x;q^{d_\ell})_{N_\ell}^{\alpha_\ell}.
\]
For \(f(x)=x(x;q)_M^\alpha\), the formula becomes
\begin{equation}\label{qcc:simple-inverse}
 [z^n]f^{-1}(z)=\frac1{n!}\Y{n-1}
   \left(\left(\alpha n(r-1)![M]_{q^r}\right)_{r=1}^{n-1}\right).
\end{equation}
An alternative finite sum, obtained by expanding each factor separately, is
\begin{equation}\label{qcc:simple-inverse-multisum}
 [z^n]f^{-1}(z)=\frac1n
 \sum_{\substack{r_0,\ldots,r_{M-1}\ge0\\\sum r_j=n-1}}
 \prod_{j=0}^{M-1}\frac{\rising{\alpha n}{r_j}}{r_j!}
 q^{\sum j r_j}.
\end{equation}
Here the rising factorial makes the formula polynomial in \(\alpha\), including
values at which gamma-quotient representations would be inconvenient.

\begin{example}[A two-color inverse]\label{qcc:inverse-example}
The inverse of \(x(1-x)(1-qx)\) begins
\begin{align*}
 y(z)={}&z+(1+q)z^2+(2+3q+2q^2)z^3\\
       &+(5+10q+10q^2+5q^3)z^4+O(z^5).
\end{align*}
At \(q=1\), the coefficient formula is
\(n^{-1}\binom{3n-2}{n-1}\), giving \(1,2,7,30,\ldots\).
\end{example}
The coefficients for positive integral \(\alpha\) count rooted plane trees
whose children are split into \(\alpha M\) ordered blocks; a child in a block of
color \(j\) has weight \(q^j\). Each vertex contributes one factor
\(\Phi(t)\), and the root contributes \(z\), proving the functional equation.
Every tree with \(n\) vertices has \(n-1\) edges. Reversing all color labels
\(j\mapsto M-1-j\) therefore proves
\begin{equation}\label{qcc:inverse-reciprocity}
 [z^n]y(z;q^{-1})=q^{-(M-1)(n-1)}[z^n]y(z;q).
\end{equation}
The same identity follows algebraically by reversing the finite product and using
homogeneity in \eqref{qcc:simple-inverse}. The tree argument additionally proves
nonnegativity and integrality in this parameter range.

\begin{auditbox}[title={Ordinary inversion is not obtained by replacing \(n\) with \([n]_q\)}]
The factor \(1/n\) in Lagrange inversion is an ordinary integer factor. For
\(M=1\), the equation is \(y=z/(1-y)\), independent of \(q\), and its
coefficients are Catalan numbers. Replacing \(1/n\) by \(1/[n]_q\) would make
the answer depend on \(q\), contradicting the defining equation. Genuine
\(q\)-shifted inversion theories change additional algebraic structure; they
are not the operation performed here.
\end{auditbox}

\section{A two-variable inverse with its Jacobian factor exposed}
The source also develops multivariate Good inversion. The following complete
specialization shows how geometric-product Bell data enter a coupled system,
while retaining the correction that a naive product of one-variable formulae
would omit.

\begin{theorem}[Cross-coupled inversion]\label{qcc:coupled}
Let \(A(0)=B(0)=1\), and let
\[
 u=z_1A(v),\qquad v=z_2B(u).
\]
For \(m,n>0\), \(r,s\ge0\), \(m\ge r\), and \(n\ge s\),
\begin{equation}\label{qcc:coupled-formula}
 [z_1^mz_2^n]u^rv^s
 =\frac{ms+nr-rs}{mn}\,
 [t^{n-s}]A(t)^m\,[t^{m-r}]B(t)^n.
\end{equation}
Each coefficient on the right has the Bell form
\eqref{qcc:master-coefficient} when \(A,B\) are finite geometric products.
\end{theorem}
\begin{proof}
Substitute the first equation into the second:
\(v=z_2B(z_1A(v))\). Apply the one-variable formula
\eqref{qcc:lagrange-formula} in \(z_2\), with \(z_1\) a parameter, to
\(u^rv^s=z_1^rA(v)^rv^s\). The derivative is
\(s t^{s-1}A(t)^r+r t^sA(t)^{r-1}A'(t)\).
Furthermore,
\[
 [z_1^{m-r}]B(z_1A(t))^n
 =A(t)^{m-r}[w^{m-r}]B(w)^n.
\]
The first derivative term therefore contributes
\(\frac{s}{n}[t^{n-s}]A(t)^m\) times the \(B\)-coefficient. The second
contributes
\[
 \frac{r}{n}[t^{n-s-1}]A(t)^{m-1}A'(t)
 =\frac{r(n-s)}{mn}[t^{n-s}]A(t)^m.
\]
Adding gives \((ms+nr-rs)/(mn)\). This argument also covers \(r=0\) or
\(s=0\), since the corresponding derivative term is zero.
\end{proof}
For comparison, the determinant in the two-variable Good formula is
\[
 1-\frac{t_2A'(t_2)}{A(t_2)}\frac{t_1B'(t_1)}{B(t_1)}.
\]
Its correction subtracts \((n-s)(m-r)/(mn)\) from one, which is exactly the
prefactor in \eqref{qcc:coupled-formula}. This gives an independent check on the
indices and on the otherwise easily omitted Jacobian contribution.

\chapter{The coefficients of a Gaussian polynomial}
\label{qcc:gaussian-coefficients}
\section{A divisor-sum formula for each coefficient}
Write \(m=n-k\), \(D=k(n-k)\), and
\[
 G_{n,k}(q)=\GB nk=\sum_{r=0}^D c_{n,k}(r)q^r.
\]
This chapter extracts coefficients in the \emph{base} \(q\), rather than in an
auxiliary product variable. Since \(G_{n,k}(0)=1\), it has a formal logarithm
in \(\Q[[q]]\).

\begin{theorem}\label{qcc:divisor-coefficients}
For \(r\ge1\), define the finite divisor sum
\begin{equation}\label{qcc:arithmetic-kernel}
 A_r(n,k)=\sum_{\substack{d\mid r\\1\le d\le k}}d
          -\sum_{\substack{d\mid r\\m<d\le n}}d.
\end{equation}
Then
\begin{equation}\label{qcc:gaussian-coeff-bell}
 c_{n,k}(r)=\frac1{r!}\Y r\left(
        \left((j-1)!A_j(n,k)\right)_{j=1}^r\right),
\end{equation}
or, with no named auxiliary polynomial,
\begin{equation}\label{qcc:gaussian-coeff-finite}
 c_{n,k}(r)=\sum_{\substack{v_1,\ldots,v_r\ge0\\\sum jv_j=r}}
                 \prod_{j=1}^r\frac{A_j(n,k)^{v_j}}{j^{v_j}v_j!}.
\end{equation}
Also \(c_{n,k}(0)=1\), and the formula gives zero for \(r>D\).
\end{theorem}
\begin{proof}
The finite product quotient gives
\begin{align*}
 \log G_{n,k}(q)
 &=\sum_{d=1}^k\sum_{a\ge1}\frac{q^{da}}a
   -\sum_{d=m+1}^n\sum_{a\ge1}\frac{q^{da}}a\\
 &=\sum_{r\ge1}\frac{A_r(n,k)}r q^r.
\end{align*}
For fixed \(r\), only divisors \(d\mid r\) contribute, with coefficient
\(1/(r/d)=d/r\). The Bell exponential formula proves the result. The original
series is a polynomial of degree \(D\), proving the last assertion.
\end{proof}
Although individual terms in \eqref{qcc:gaussian-coeff-finite} are rational and
can have signs, the answer is a nonnegative integer by the subset definition.
For example, when \((n,k)=(4,2)\), the first kernels are
\(A_1=1,A_2=3,A_3=-2,A_4=-1\). They yield
\[
 c(1)=1,\qquad c(2)=\frac{1+3}{2}=2,\qquad
 c(3)=\frac{1+9-4}{6}=1,
\]
and \(c(4)=1\), as expected.

\section{A second finite formula by inclusion--exclusion}
\begin{proposition}\label{qcc:coefficient-inclusion}
For \(r\ge0\),
\begin{equation}\label{qcc:coefficient-inclusion-formula}
 c_{n,k}(r)=\sum_{\varepsilon\in\{0,1\}^k}(-1)^{\sum\varepsilon_j}
 \sum_{\substack{v_1,\ldots,v_k\ge0\\
       \sum jv_j=r-\sum(m+j)\varepsilon_j}}1.
\end{equation}
An impossible inner constraint contributes zero. All inner variables may be
bounded by \(v_j\le\lfloor r/j\rfloor\).
\end{proposition}
\begin{proof}
Expand each numerator factor \(1-q^{m+j}\) by its two terms and each reciprocal
factor \((1-q^j)^{-1}\) as a geometric series. The sign records which numerator
terms were selected, and the displayed constraint is precisely the total
exponent. Coefficient extraction proves the formula.
\end{proof}
The two finite formulae organize the same coefficient differently. The first
collects repeated divisibility contributions into Bell blocks; the second uses
restricted partitions and inclusion--exclusion. Neither needs earlier values of
\(c_{n,k}(r)\).

\section{Finite Fourier extraction}
Let \(L>D\), and let \(\omega=e^{2\pi i/L}\). Orthogonality of roots of unity
implies, for \(0\le r\le D\),
\begin{equation}\label{qcc:fourier-coefficient}
 c_{n,k}(r)=\frac1L\sum_{j=0}^{L-1}\omega^{-rj}G_{n,k}(\omega^j).
\end{equation}
Indeed, expand the polynomial and use
\(L^{-1}\sum_{j=0}^{L-1}\omega^{j(s-r)}=\delta_{sr}\) in this degree range.
Without the condition \(L>D\), the same expression gives the sum of coefficients
with exponent congruent to \(r\pmod L\). Later, \(q\)-Lucas supplies a direct
root-of-unity evaluation of every summand.

This identity is exact in a cyclotomic field. Floating-point use can suffer from
cancellation, so it is not automatically a preferable numerical algorithm. For
exact coefficient production, polynomial Pascal recursion or the divisor-sum
formula is usually simpler.

\chapter{Parameter derivatives and the expansion at one}
\label{qcc:parameter-jets}
\section{Why the logarithmic parameter is the natural coordinate}
Put \(m=n-k\), \(D=km\), and write
\[
 G_{n,k}(q)=\GB nk=\sum_{j=0}^{D}c_{n,k}(j)q^j.
\]
The substitution \(q=e^t\) converts powers of \(q\) into exponentials. Thus
ordinary derivatives in \(t\) give power moments of the coefficient sequence,
whereas derivatives in \(q\) give falling-factorial moments. This is precisely
the distinction for which the source article uses the Stirling basis changes.
Let \(\logq=q\,\mathrm d/\mathrm dq\), so that
\(\frac{\mathrm d}{\mathrm dt}F(e^t)=(\logq F)(e^t)\).

\begin{lemma}[The regular logarithm of an exponential difference]
\label{qcc:regular-log}
As a formal series at \(u=0\),
\begin{equation}\label{qcc:regular-log-eq}
 \log\frac{e^u-1}{u}
 =\frac u2+\sum_{r\ge2}\frac{\beta_r}{r}\frac{u^r}{r!}.
\end{equation}
All odd coefficients after the linear one vanish.
\end{lemma}
\begin{proof}
The quotient inside the logarithm has constant term one. Its logarithmic
 derivative is
\[
 \frac{e^u}{e^u-1}-\frac1u
 =1+\frac1u\left(\frac{u}{e^u-1}-1\right)
 =\frac12+\sum_{r\ge2}\beta_r\frac{u^{r-1}}{r!}.
\]
Integrate formally and use the zero constant term of the logarithm. The function
\(u/(e^u-1)+u/2\) is even, as direct substitution of \(-u\) shows. Hence
\(\beta_{2j+1}=0\) for \(j\ge1\).
\end{proof}

\begin{theorem}[All-order Gaussian cumulants]
\label{qcc:cumulants}
Define
\begin{equation}\label{qcc:kappa}
 \kappa_1(n,k)=\frac{km}{2},\qquad
 \kappa_r(n,k)=\frac{\beta_r}{r}
       \sum_{j=1}^{k}\bigl((m+j)^r-j^r\bigr)\quad(r\ge2).
\end{equation}
Then
\begin{align}
 \log\frac{G_{n,k}(e^t)}{\binom nk}
   &=\sum_{r\ge1}\kappa_r(n,k)\frac{t^r}{r!},
       \label{qcc:log-cumulants}\\
 G_{n,k}(e^t)
   &=\binom nk\sum_{r\ge0}
      \Y r(\kappa_1,\ldots,\kappa_r)\frac{t^r}{r!}.
       \label{qcc:cumulant-expansion}
\end{align}
These are both formal identities and convergent local Taylor expansions.
In particular,
\begin{equation}\label{qcc:low-cumulants}
 \kappa_2=\frac{D(n+1)}{12},\qquad
 \kappa_3=0,\qquad
 \kappa_4=-\frac{D(n+1)\bigl(n(n+1)-D\bigr)}{120}.
\end{equation}
\end{theorem}
\begin{proof}
The cancelled product representation is
\[
 \frac{G_{n,k}(e^t)}{\binom nk}
 =\prod_{j=1}^{k}
 \frac{(e^{(m+j)t}-1)/((m+j)t)}{(e^{jt}-1)/(jt)}.
\]
Every factor now has constant term one. Apply
\cref{qcc:regular-log} to each numerator and denominator. The linear
coefficient is \(\frac12\sum_{j=1}^k m=km/2\); the higher coefficients
are exactly \eqref{qcc:kappa}. Exponentiation and the Bell generating
function give \eqref{qcc:cumulant-expansion}.

For the second coefficient,
\(\sum_{j=1}^k((m+j)^2-j^2)=km(n+1)\). For the fourth,
expansion of the fourth powers and summation of powers of \(j\) gives
\[
 \sum_{j=1}^k((m+j)^4-j^4)
 =km(n+1)\bigl(n(n+1)-km\bigr).
\]
Together with \(\beta_2=1/6\), \(\beta_3=0\), and \(\beta_4=-1/30\),
these are \eqref{qcc:low-cumulants}. The local analytic assertions follow
because the cancelled product is analytic and nonzero near \(t=0\).
\end{proof}

\subsection{Finite Bernoulli-polynomial form}
No power-sum recurrence is necessary in \eqref{qcc:kappa}. Define
\(\beta_r(x)=\sum_{j=0}^r\binom rj\beta_jx^{r-j}\). Multiplying its
generating function \(te^{xt}/(e^t-1)\) by \(e^t-1\) shows
\(\beta_{r+1}(x+1)-\beta_{r+1}(x)=(r+1)x^r\). Consequently
\[
 P_r(M):=\sum_{j=1}^{M}j^r
   =\frac{\beta_{r+1}(M+1)-\beta_{r+1}(1)}{r+1}.
\]
The sum in \eqref{qcc:kappa} is \(P_r(n)-P_r(m)-P_r(k)\).
This also makes its symmetry in \(k,m\) manifest.

\section{Exact coefficient moments and probability}
\begin{corollary}[Power and factorial moments]\label{qcc:moment-formulas}
For \(r\ge0\),
\begin{align}
 \sum_{j=0}^{D}j^r c_{n,k}(j)
 &=\binom nk\Y r(\kappa_1,\ldots,\kappa_r),
       \label{qcc:power-moments}\\
 G_{n,k}^{(r)}(1)
 &=\binom nk\sum_{a=0}^r s(r,a)
           \Y a(\kappa_1,\ldots,\kappa_a).
       \label{qcc:ordinary-jets-one}
\end{align}
For \(r=0\), the first sum is the coefficient sum, with \(0^0=1\).
\end{corollary}
\begin{proof}
Differentiate the finite sum \(G_{n,k}(e^t)=\sum_jc_{n,k}(j)e^{jt}\).
This proves \eqref{qcc:power-moments}. The identity
\(j^{\underline r}=\sum_a s(r,a)j^a\), which defines the signed
Stirling coefficients, gives \eqref{qcc:ordinary-jets-one}.
\end{proof}

Choosing a partition uniformly from the partitions in the \(k\)-by-\(m\)
rectangle gives a random area \(X\) with
\(\Pr(X=j)=c_{n,k}(j)/\binom nk\). Its moment generating function is the
left side of \eqref{qcc:cumulant-expansion}, divided by \(\binom nk\).
Thus the \(\kappa_r\) are its cumulants, and
\begin{equation}\label{qcc:area-statistics}
 \E X=\frac D2,\qquad \Var(X)=\frac{D(n+1)}{12}.
\end{equation}
Complementing the rectangle sends \(X\) to \(D-X\); hence every odd
central moment is zero. Every even central moment is given by
\(\Y{2r}(0,\kappa_2,0,\kappa_4,\ldots,\kappa_{2r})\).
These assertions concern a finite distribution and require no limiting
probabilistic theorem.

\begin{example}
For \(G_{4,2}(q)=1+q+2q^2+q^3+q^4\), the coefficient sum is six,
\(\kappa_1=2\), \(\kappa_2=5/3\), and \(\kappa_4=-8/3\).
Therefore
\[
 G_{4,2}'(1)=6\cdot2=12,\qquad
 (\logq^2G_{4,2})(1)=6(2^2+5/3)=34,
\]
whereas
\(G_{4,2}''(1)=34-12=22\). Confusing the two kinds of differentiation
would already give an incorrect second derivative.
\end{example}

\section{Arbitrary parameter points and Eulerian rational functions}
For \(s\ge0\), define the rational function
\[
 \Li_{-s}(u)=\left(u\frac{\mathrm d}{\mathrm du}\right)^s\frac{u}{1-u}.
\]
For \(|u|<1\) this is \(\sum_{j\ge1}j^su^j\), but the rational
expression is the definition needed outside that disk. Set \(A_0(u)=1\).
Then
\begin{equation}\label{qcc:li-eulerian}
 \Li_{-s}(u)=\frac{uA_s(u)}{(1-u)^{s+1}},\qquad
 A_{s+1}(u)=(1+su)A_s(u)+u(1-u)A_s'(u).
\end{equation}
The formula follows by applying \(u\,\mathrm d/\mathrm du\) to the
fraction. For \(s\ge1\), the coefficients of \(A_s\) are the ordinary
Eulerian numbers indexed by the number of descents: insertion of the largest
letter gives the same coefficient recurrence and initial value. Alternatively,
\eqref{qcc:li-eulerian} is an independent algebraic definition of this
Eulerian normalization.

\begin{theorem}[Logarithmic parameter derivatives]\label{qcc:generic-jets}
At a nonzero point \(q\) for which none of \(q^1,\ldots,q^n\) equals
one, define
\begin{equation}\label{qcc:L-r}
 L_r(q)=\sum_{j=1}^{k}
 \left(j^r\Li_{1-r}(q^j)-(m+j)^r\Li_{1-r}(q^{m+j})\right).
\end{equation}
Then, for \(r\ge0\),
\begin{align}
 \logq^rG_{n,k}(q)&=G_{n,k}(q)\Y r(L_1(q),\ldots,L_r(q)),
          \label{qcc:theta-jets}\\
 G_{n,k}^{(r)}(q)&=q^{-r}G_{n,k}(q)
       \sum_{a=0}^r s(r,a)\Y a(L_1(q),\ldots,L_a(q)).
          \label{qcc:q-jets}
\end{align}
All \(L_r\) are explicit rational functions, by
\eqref{qcc:li-eulerian}.
\end{theorem}
\begin{proof}
For \(a\ge1\), repeated logarithmic differentiation gives
\[
 \logq^r\log(1-q^a)=-a^r\Li_{1-r}(q^a),\qquad r\ge1.
\]
Subtracting denominator contributions from numerator contributions gives
\(\logq^r\log G=L_r\). In the local coordinate \(q e^t\), the
logarithm of \(G(qe^t)/G(q)\) is
\(\sum_{r\ge1}L_r(q)t^r/r!\). Exponentiate and use Bell polynomials.
Finally, the operator identity
\begin{equation}\label{qcc:stirling-operators}
 q^r\frac{\mathrm d^r}{\mathrm dq^r}
 =\logq(\logq-1)\cdots(\logq-r+1)
 =\sum_{a=0}^r s(r,a)\logq^a
\end{equation}
follows by applying both sides to each monomial. It proves the second formula.
\end{proof}

The displayed domain avoids unnecessary cancellation inside individual rational
terms. At roots of unity, the Gaussian polynomial is still regular, but some
\(L_r\) may have poles. The correct replacement is a deflated local
factorization, not termwise substitution into \eqref{qcc:L-r}.

\chapter{Cyclotomic structure and all-order root-of-unity jets}
\label{qcc:roots}
\section{Cyclotomic factorization and carries}
Let \(\Phi_d(q)\) denote the monic polynomial whose roots are the primitive
\(d\)-th roots of unity. The partition of roots by their exact order gives
\(q^a-1=\prod_{d\mid a}\Phi_d(q)\).

\begin{theorem}[Squarefree cyclotomic factorization]\label{qcc:cyclotomic}
For \(0\le k\le n\), \(m=n-k\),
\begin{equation}\label{qcc:cyclotomic-product}
 G_{n,k}(q)=\prod_{d=2}^{n}\Phi_d(q)^{\epsilon_d},\qquad
 \epsilon_d=\left\lfloor\frac nd\right\rfloor
      -\left\lfloor\frac kd\right\rfloor
      -\left\lfloor\frac md\right\rfloor\in\{0,1\}.
\end{equation}
Consequently every nonconstant Gaussian polynomial has distinct roots, all on
the unit circle. Its degree gives the identity
\begin{equation}\label{qcc:totient-degree}
 \sum_{d=2}^{n}\varphi(d)\epsilon_d=k(n-k).
\end{equation}
\end{theorem}
\begin{proof}
Factor each \(1-q^a\) in the factorial quotient. There are
\(\lfloor n/d\rfloor\) multiples of \(d\) in its numerator factorial,
and the denominator subtracts the corresponding counts for \(k\) and \(m\).
The exponent of \(\Phi_1=q-1\) is \(n-k-m=0\); the signs also cancel.
Write \(k=ud+v\), \(m=wd+x\), with \(0\le v,x<d\). Then
\(\epsilon_d=\lfloor(v+x)/d\rfloor\), which is either zero or one.
The roots of different cyclotomic polynomials are disjoint, and each is simple.
Finally \(\deg\Phi_d=\varphi(d)\), and degrees add in the product.
\end{proof}
This carry interpretation is particularly useful: vanishing at a primitive
\(d\)-th root occurs precisely when the residues of \(k\) and \(n-k\)
produce a carry upon addition modulo \(d\).

\section{The finite \qbname-Lucas theorem}
The following classical root-of-unity formula is often called \(q\)-Lucas;
see, for example, \cite{formichella-straub}. Its proof is short enough to
include, and the phase bookkeeping is important.

\begin{theorem}[\(q\)-Lucas]\label{qcc:lucas}
Let \(\zeta\) be a primitive \(d\)-th root of unity, with \(d\ge1\).
Write \(n=ad+b\), \(k=rd+s\), where \(0\le b,s<d\). Then
\begin{equation}\label{qcc:lucas-eq}
 G_{n,k}(\zeta)=\binom ar\,G_{b,s}(\zeta).
\end{equation}
The right side is zero when \(s>b\).
\end{theorem}
\begin{proof}
The finite product theorem and a division into full blocks give
\[
 \prod_{j=0}^{n-1}(1+\zeta^jz)
 =\bigl(1-(-z)^d\bigr)^a
       \prod_{j=0}^{b-1}(1+\zeta^jz).
\]
The second factor has degree smaller than \(d\). Hence its only possible
contribution to \(z^{rd+s}\) has degree \(s\), and coefficient comparison
therefore gives
\[
 \zeta^{\binom k2}G_{n,k}(\zeta)
 =(-1)^{r(d+1)}\binom ar
           \zeta^{\binom s2}G_{b,s}(\zeta).
\]
The phase identity
\(\zeta^{\binom{rd+s}{2}-\binom s2}=(-1)^{r(d+1)}\)
finishes the proof. For odd \(d\), the exponent is a multiple of \(d\).
For even \(d\), it is congruent to \(rd/2\) modulo \(d\), giving
\((-1)^r\). These are exactly the two values on the right.
The same argument covers \(d=1\).
\end{proof}

For example,
\[
 G_{2a,2r}(-1)=\binom ar,\qquad G_{2a,2r+1}(-1)=0,
\]
while \(G_{2a+1,2r}(-1)=G_{2a+1,2r+1}(-1)=\binom ar\).
If \(d\mid n\), \eqref{qcc:lucas-eq} is also a fixed-subset count: a
permutation consisting of \(n/d\) disjoint \(d\)-cycles fixes a
\(k\)-element subset exactly when the subset is a union of whole cycles.
There are \(\binom{n/d}{k/d}\) such subsets if \(d\mid k\), and none
otherwise. This provides a concrete combinatorial interpretation for these
root evaluations.

\section{Deflating all apparent singularities}
A value formula such as \(q\)-Lucas does not determine derivatives at the
root. The next theorem supplies every derivative, including the case where the
Gaussian polynomial itself vanishes. The calculation uses the logarithmic
coordinate \(q=\zeta e^t\) and removes powers of \(t\) before taking a
logarithm.

For an integer \(a\ge1\), define
\begin{equation}\label{qcc:root-constants}
 c_a(\zeta)=
 \begin{cases}
 -a,&d\mid a,\\
 1-\zeta^a,&d\nmid a.
 \end{cases}
\end{equation}
For \(r\ge1\), define
\begin{equation}\label{qcc:eta}
 \eta_r(a;\zeta)=
 \begin{cases}
 a/2,&d\mid a,\ r=1,\\[1mm]
 \beta_r a^r/r,&d\mid a,\ r\ge2,\\[1mm]
 -a^r\Li_{1-r}(\zeta^a),&d\nmid a.
 \end{cases}
\end{equation}
All rational functions in the last line are evaluated away from their pole at
one. Put
\begin{align}
 \nu&=\left\lfloor\frac nd\right\rfloor
       -\left\lfloor\frac kd\right\rfloor
       -\left\lfloor\frac md\right\rfloor,\label{qcc:root-order}\\
 C_\zeta&=\prod_{j=1}^{k}\frac{c_{m+j}(\zeta)}{c_j(\zeta)},\label{qcc:root-C}\\
 \lambda_r&=\sum_{j=1}^{k}
       \bigl(\eta_r(m+j;\zeta)-\eta_r(j;\zeta)\bigr).
       \label{qcc:root-lambda}
\end{align}
Here \(\nu=0\) when \(d=1\), and always \(C_\zeta\ne0\).

\begin{theorem}[Uniform all-order local formula]\label{qcc:root-jets}
With the definitions above,
\begin{equation}\label{qcc:root-factorization}
 \boxed{\quad
 G_{n,k}(\zeta e^t)
  =C_\zeta t^\nu
     \exp\left(\sum_{r\ge1}\lambda_r\frac{t^r}{r!}\right).
 \quad}
\end{equation}
In particular, for \(r\ge0\),
\begin{equation}\label{qcc:root-t-coefficients}
 [t^{\nu+r}]G_{n,k}(\zeta e^t)
       =\frac{C_\zeta}{r!}\Y r(\lambda_1,\ldots,\lambda_r).
\end{equation}
Every ordinary Taylor derivative at \(\zeta\) is the finite sum
\begin{equation}\label{qcc:root-ordinary-derivatives}
 G_{n,k}^{(R)}(\zeta)
 =C_\zeta\zeta^{-R}
   \sum_{j=\nu}^{R}s(R,j)\frac{j!}{(j-\nu)!}
               \Y{j-\nu}(\lambda_1,\ldots,\lambda_{j-\nu}).
\end{equation}
An empty sum is zero. Formula \eqref{qcc:root-factorization} holds formally
in \(\Q(\zeta)[[t]]\) and analytically near \(t=0\).
\end{theorem}
\begin{proof}
For each factor in the product quotient, consider two cases. If \(d\mid a\),
then
\[
 1-\zeta^ae^{at}=-at\,\frac{e^{at}-1}{at}
 =c_a(\zeta)t
     \exp\left(\frac{at}{2}
          +\sum_{r\ge2}\frac{\beta_ra^r}{r}\frac{t^r}{r!}\right)
\]
by \cref{qcc:regular-log}. If \(d\nmid a\), the factor has nonzero
constant term \(c_a(\zeta)\), and repeated differentiation yields
\[
 \log\frac{1-\zeta^ae^{at}}{1-\zeta^a}
 =-\sum_{r\ge1}a^r\Li_{1-r}(\zeta^a)\frac{t^r}{r!}.
\]
Multiplying numerator factors and dividing denominator factors produces the
constant \(C_\zeta\), the net power \(t^\nu\), and the sum of logarithms
\(\sum\lambda_rt^r/r!\). This proves \eqref{qcc:root-factorization};
Bell expansion proves \eqref{qcc:root-t-coefficients}.

The operator \(\logq\) becomes \(\mathrm d/\mathrm dt\). Therefore
\[
 (\logq^jG)(\zeta)=
 \begin{cases}
 0,&j<\nu,\\
 C_\zeta\dfrac{j!}{(j-\nu)!}
        \Y{j-\nu}(\lambda_1,\ldots,\lambda_{j-\nu}),&j\ge\nu.
 \end{cases}
\]
Apply \eqref{qcc:stirling-operators} at \(q=\zeta\) to get
\eqref{qcc:root-ordinary-derivatives}. Finally, after deflation all factors
are analytic units near zero, so their normalized logarithms converge locally.
No logarithm of a vanishing quantity was used.
\end{proof}

\begin{corollary}[First two derivatives]\label{qcc:root-low-jets}
At a nonzero value, \(\nu=0\),
\[
 G(\zeta)=C_\zeta,\qquad
 G'(\zeta)=C_\zeta\zeta^{-1}\lambda_1,\qquad
 G''(\zeta)=C_\zeta\zeta^{-2}
             (\lambda_1^2+\lambda_2-\lambda_1).
\]
At a zero, necessarily simple, \(\nu=1\),
\[
 G'(\zeta)=C_\zeta\zeta^{-1},\qquad
 G''(\zeta)=C_\zeta\zeta^{-2}(2\lambda_1-1).
\]
\end{corollary}
\begin{proof}
Substitute \(R=0,1,2\) into
\eqref{qcc:root-ordinary-derivatives}, using
\(s(2,1)=-1\), \(s(2,2)=1\), and
\(\Y2(x_1,x_2)=x_1^2+x_2\).
\end{proof}

\begin{example}[A zero at \(-1\)]
Take \(G_{4,1}(q)=1+q+q^2+q^3\). At \(\zeta=-1\),
\[
 \nu=1,\qquad C_{-1}=-2,\qquad
 \lambda_1=\frac32,\qquad\lambda_2=\frac{13}{12}.
\]
Thus \(G'(-1)=2\) and \(G''(-1)=-4\), in agreement with
\[
 G_{4,1}(-1+h)=2h-2h^2+h^3.
\]
\end{example}

At \(\zeta=1\), \(C_1=\binom nk\), \(\nu=0\), and
\(\lambda_r=\kappa_r(n,k)\): the expansion at one is recovered. Together,
\cref{qcc:gaussian-coefficients,qcc:generic-jets,qcc:root-jets} give explicit
coefficients at every complex base.

\chapter{Weighted sums and differentiated identities}
\label{qcc:weighted-sums}
\section{Polynomial weights in the summation index}
The finite product theorem can be differentiated in its auxiliary variable
without differentiating the parameter \(q\). This gives a second, independent
source of weighted Gaussian sums. Let
\[
 P_N(z)=(-z;q)_N
   =\sum_{k=0}^{N}q^{\binom k2}\GB Nk z^k,
 \qquad \Theta_z=z\frac{\mathrm d}{\mathrm dz}.
\]

\begin{theorem}[All power-weighted finite sums]\label{qcc:weighted-Gaussian}
At a point where \(P_N(z)\ne0\), define
\[
 K_r(z)=-\sum_{j=0}^{N-1}\Li_{1-r}(-zq^j),\qquad r\ge1.
\]
Then, for \(r\ge0\),
\begin{equation}\label{qcc:weighted-Gaussian-eq}
 \sum_{k=0}^{N}k^r q^{\binom k2}\GB Nk z^k
   =P_N(z)\Y r(K_1(z),\ldots,K_r(z)).
\end{equation}
The right side, after cancellation, is a polynomial and therefore extends to
the zeros of \(P_N\).
\end{theorem}
\begin{proof}
The left side is \(\Theta_z^rP_N(z)\). For every \(r\ge1\),
\(\Theta_z^r\log P_N(z)=K_r(z)\). Taylor expansion of
\(P_N(ze^t)/P_N(z)\) and the Bell generating function give the result.
The polynomial extension follows because the two rational expressions agree
away from finitely many zeros.
\end{proof}
For example,
\[
 \sum_{k=0}^{N}kq^{\binom k2}\GB Nk z^k
  =(-z;q)_N\sum_{j=0}^{N-1}\frac{zq^j}{1+zq^j}.
\]
For a general polynomial \(p(k)=\sum_{r=0}^{R}a_rk^r\), sum
\eqref{qcc:weighted-Gaussian-eq} with weights \(a_r\). Alternatively,
expand \(p\) in falling factorials and use ordinary derivatives in \(z\).

\section{Alternating sums at a product zero}
At \(z=1\), the product \((z;q)_N\) vanishes for \(N\ge1\), so a
logarithmic formula should first be deflated. Set
\[
 H_N(z)=(zq;q)_{N-1},\qquad (z;q)_N=(1-z)H_N(z).
\]
Assume for the moment that \(q^j\ne1\) for \(1\le j<N\), and put
\begin{equation}\label{qcc:alternating-kernels}
 J_s=-(s-1)!\sum_{j=1}^{N-1}
             \frac{q^{js}}{(1-q^j)^s}\quad(s\ge1).
\end{equation}

\begin{proposition}[Alternating factorial moments]\label{qcc:alternating}
For \(r\ge1\),
\begin{equation}\label{qcc:alternating-eq}
 \sum_{k=0}^{N}(-1)^kq^{\binom k2}\GB Nk\fall{k}{r}
   =-r(q;q)_{N-1}\Y{r-1}(J_1,\ldots,J_{r-1}).
\end{equation}
For \(r=0\), the sum is zero. Ordinary powers are recovered by
\(k^r=\sum_{a=0}^{r}S(r,a)\fall{k}{a}\).
\end{proposition}
\begin{proof}
The left side is the \(r\)-th derivative of \((z;q)_N\) at one.
Leibniz's rule and the factor \(1-z\) leave only
\(-rH_N^{(r-1)}(1)\). The ordinary derivatives of
\(\log H_N(z)\) at one are exactly \(J_s\), and
\(H_N(1)=(q;q)_{N-1}\). Apply the ordinary Bell derivative formula,
which follows from the Taylor expansion of \(H_N(1+t)/H_N(1)\).
The zeroth-order sum is \((1;q)_N=0\). The conversion of powers to
falling factorials is the defining ordinary second-kind Stirling expansion.
\end{proof}

When \(r>N\), the left side is zero, and the right side is also zero:
it represents a derivative of the degree-\(N-1\) polynomial \(H_N\)
of order \(r-1>N-1\). At excluded roots of unity, polynomial
specialization or the root deflation of \cref{qcc:root-jets} replaces
termwise evaluation of \eqref{qcc:alternating-kernels}.

\begin{auditbox}[title={Which cancellation is actually present?}]
The Gaussian alternating kernel does not simply annihilate every polynomial
weight of degree less than \(N\). Already at \(N=2\),
\[
 \sum_{k=0}^{2}(-1)^kq^{\binom k2}\GB2k k=q-1,
\]
which is nonzero for generic \(q\). What the product theorem gives is
cancellation at specified \emph{geometric spectral points}. The Jackson stencil
and the extrapolation filters below include the necessary additional powers of
\(q\). Replacing every ordinary binomial coefficient in an additive
finite-difference identity by a Gaussian coefficient is not a valid rule.
\end{auditbox}

\section{An all-order differentiation of central Vandermonde}
The central convolution proved earlier is
\[
 G_{2n,n}(q)=\sum_{k=0}^{n}q^{k^2}G_{n,k}(q)^2.
\]
Differentiating this identity at \(q=1\) produces infinitely many finite
identities involving ordinary binomial coefficients, Bernoulli numbers, and
Bell polynomials. The logarithmic coordinate makes the formula uniform.

\begin{theorem}[Differentiated central convolution]
\label{qcc:differentiated-central}
For \(r\ge0\), define, for each \(k\),
\[
 v_1^{(k)}=nk,\qquad
 v_a^{(k)}=2\kappa_a(n,k)\quad(a\ge2).
\]
Then
\begin{equation}\label{qcc:differentiated-central-eq}
 \binom{2n}{n}\Y r\bigl(\kappa_1(2n,n),\ldots,\kappa_r(2n,n)\bigr)
 =\sum_{k=0}^{n}\binom nk^2
           \Y r\bigl(v_1^{(k)},\ldots,v_r^{(k)}\bigr).
\end{equation}
\end{theorem}
\begin{proof}
Substitute \(q=e^t\). On the \(k\)-th summand the logarithm, after
removing the constant \(\binom nk^2\), is
\[
 k^2t+2\sum_{a\ge1}\kappa_a(n,k)\frac{t^a}{a!}.
\]
Its linear coefficient is
\(k^2+2\kappa_1(n,k)=k^2+k(n-k)=nk\), and its remaining cumulants
are \(2\kappa_a(n,k)\). Expand both sides by
\cref{qcc:cumulants} and compare coefficients of \(t^r/r!\).
\end{proof}

For \(r=1\) and \(n\ge1\), this gives
\[
 \sum_{k=0}^{n} k\binom nk^2=\frac n2\binom{2n}{n}.
\]
For \(r=2\), it gives the less immediate identity
\[
 \sum_{k=0}^{n}\binom nk^2
    \left(n^2k^2+\frac{k(n-k)(n+1)}6\right)
 =\binom{2n}{n}\left(\frac{n^4}{4}
                      +\frac{n^2(2n+1)}{12}\right).
\]
Higher identities require no new differentiation argument: the same finite Bell
formula gives their complete coefficients. The same procedure works for every
block convolution in \cref{qcc:convolution}; each explicit power of \(q\)
adds only to the first cumulant.

\chapter{Gaussian filters on geometric grids}
\label{qcc:filters}
\section{Spectral design instead of recursive extrapolation}
Let \(0<\rho<1\) and define the dilation operator
\((T_\rho f)(h)=f(\rho h)\). Its eigenfunctions are the monomials:
\(T_\rho h^m=\rho^m h^m\). To preserve a constant and annihilate the
first \(N\) positive powers, construct the normalized spectral polynomial
\begin{equation}\label{qcc:filter-def}
 R_N=\prod_{a=1}^{N}\frac{T_\rho-\rho^a}{1-\rho^a}.
\end{equation}
Here \(R_0\) is the identity operator. This is a geometric version of the
source's principle of choosing a basis in which an operator is diagonal.

\begin{theorem}[Closed Gaussian extrapolation weights]
\label{qcc:filter-weights}
For every \(N\ge0\),
\begin{align}
 (R_Nf)(h)&=\sum_{j=0}^{N}w_{N,j}f(\rho^jh),
          \label{qcc:filter-sum}\\
 w_{N,j}&=
 \frac{(-1)^{N-j}\rho^{\binom{N-j+1}{2}}
                \GB[\rho]Nj}{(\rho;\rho)_N}.
          \label{qcc:weights}
\end{align}
The filter preserves constants and annihilates \(h^m\) for
\(1\le m\le N\). For every integer \(m\ge N+1\), its complete
remaining response is
\begin{equation}\label{qcc:filter-response}
 R_Nh^m=(-1)^N\rho^{\binom{N+1}{2}}
              \GB[\rho]{m-1}{N}h^m.
\end{equation}
\end{theorem}
\begin{proof}
Expand the numerator polynomial in \eqref{qcc:filter-def}. The coefficient
of \(T_\rho^j\) is
\((-1)^{N-j}e_{N-j}(\rho,\ldots,\rho^N)\). Principal specialization
of the elementary symmetric function gives
\[
 e_{N-j}(\rho,\ldots,\rho^N)
 =\rho^{\binom{N-j+1}{2}}\GB[\rho]{N}{N-j},
\]
and Gaussian symmetry gives \eqref{qcc:weights}. Evaluate the spectral
polynomial at \(1,\rho,\ldots,\rho^N\) to prove preservation and
annihilation. For \(m\ge N+1\), direct factorization gives
\[
 \prod_{a=1}^{N}\frac{\rho^m-\rho^a}{1-\rho^a}
 =(-1)^N\rho^{N(N+1)/2}
      \prod_{a=1}^{N}\frac{1-\rho^{m-a}}{1-\rho^a}
 =(-1)^N\rho^{\binom{N+1}{2}}\GB[\rho]{m-1}{N}.
\]
This is \eqref{qcc:filter-response}.
\end{proof}

\begin{corollary}[Exact analytic error series and a bound]
\label{qcc:filter-error}
Suppose \(f(h)=L+\sum_{m\ge1}c_mh^m\) is analytic for \(|h|<R\).
Then, throughout that disk,
\begin{equation}\label{qcc:filter-error-series}
 (R_Nf)(h)-L=(-1)^N\rho^{\binom{N+1}{2}}
       \sum_{m\ge N+1}\GB[\rho]{m-1}{N}c_mh^m.
\end{equation}
If, additionally, \(f\) is analytic on a neighborhood of the closed disk
\(|h|\le R\) and \(|f(h)|\le M\) there, then, for \(x=|h|/R<1\),
\begin{equation}\label{qcc:filter-error-bound}
 |(R_Nf)(h)-L|
 \le M\rho^{\binom{N+1}{2}}
                  \frac{x^{N+1}}{(x;\rho)_{N+1}}.
\end{equation}
\end{corollary}
\begin{proof}
The operator is a finite combination of evaluations at points inside the disk,
so it can be applied term by term to the convergent power series. Use
\eqref{qcc:filter-response}. Cauchy's coefficient bound gives
\(|c_m|\le M R^{-m}\). All Gaussian values are nonnegative at
\(0<\rho<1\), and the reciprocal finite-product theorem gives
\[
 \sum_{m\ge N+1}\GB[\rho]{m-1}{N}x^m
   =\frac{x^{N+1}}{(x;\rho)_{N+1}}.
\]
Taking absolute values therefore proves \eqref{qcc:filter-error-bound}.
\end{proof}

\section{Noise amplification is bounded at a fixed geometric ratio}
\begin{theorem}[Exact absolute-weight norm]\label{qcc:filter-noise}
For \(0<\rho<1\),
\begin{equation}\label{qcc:noise-norm}
 \sum_{j=0}^{N}|w_{N,j}|
       =\frac{(-\rho;\rho)_N}{(\rho;\rho)_N}
       \le\prod_{a\ge1}\frac{1+\rho^a}{1-\rho^a}<\infty.
\end{equation}
Thus the bound is uniform in \(N\) when \(\rho\) is fixed. One explicit,
though not sharp, estimate is
\begin{equation}\label{qcc:noise-simple-bound}
 \sum_{j=0}^{N}|w_{N,j}|
       \le\exp\left(\frac{2\rho}{(1-\rho)^2}\right).
\end{equation}
If each sample has absolute error at most \(\varepsilon\), the resulting
filtered error is at most \(\varepsilon\) times
\eqref{qcc:noise-norm}.
\end{theorem}
\begin{proof}
After taking absolute values, the numerator sum in \eqref{qcc:weights} is
\[
 \sum_{b=0}^{N}\rho^{\binom{b+1}{2}}\GB[\rho]Nb
       =(-\rho;\rho)_N
\]
by the finite product theorem. For \(0\le u<1\),
\[
 \log\frac{1+u}{1-u}
 =2\sum_{r\ge0}\frac{u^{2r+1}}{2r+1}
 \le\frac{2u}{1-u}.
\]
Consequently the logarithm of the infinite product is at most
\(\sum_{a\ge1}2\rho^a/(1-\rho^a)
 \le2\rho/(1-\rho)^2\), which proves convergence and the bound.
The final assertion is the triangle inequality applied to
\eqref{qcc:filter-sum}.
\end{proof}

The fixed-ratio hypothesis matters. The bound need not remain moderate as
\(\rho\) approaches one. Nor does it control the error in generating the
samples themselves; the theorem assumes a bound on those errors and describes
only the linear combination's amplification.

\begin{example}[A dyadic three-point formula]
For \(\rho=1/2\) and \(N=2\),
\begin{equation}\label{qcc:dyadic-filter}
 (R_2f)(h)=\frac13 f(h)-2f(h/2)+\frac83f(h/4).
\end{equation}
It preserves constants, cancels the linear and quadratic terms, and has leading
error \(c_3h^3/8\). Its exact absolute-weight norm is five. More generally,
\eqref{qcc:filter-error-series} specifies \emph{every} remaining power,
not just the leading one.
\end{example}

\section{Fractional powers and arithmetic exponent lattices}
For positive \(h\), let the unwanted exponents be
\(\alpha,\alpha+\delta,\ldots,\alpha+(N-1)\delta\), where
\(\alpha,\delta>0\). Define
\begin{equation}\label{qcc:general-filter}
 R_{N;\alpha,\delta}
 =\prod_{a=0}^{N-1}
       \frac{T_\rho-\rho^{\alpha+a\delta}}
            {1-\rho^{\alpha+a\delta}}.
\end{equation}
It preserves constants and annihilates exactly those prescribed power modes.
The same elementary-symmetric specialization gives the closed weights
\begin{equation}\label{qcc:general-weights}
 [T_\rho^j]R_{N;\alpha,\delta}
 =\frac{(-1)^{N-j}
        \rho^{\alpha(N-j)+\delta\binom{N-j}{2}}
        \GB[\rho^\delta]Nj}
       {(\rho^\alpha;\rho^\delta)_N}.
\end{equation}
This is proved exactly as \eqref{qcc:weights}, now with the alphabet
\(\rho^\alpha,\rho^{\alpha+\delta},\ldots\). For a general exponent
\(\beta\), the full response is the explicit product
\[
 R_{N;\alpha,\delta}h^\beta
  =h^\beta\prod_{a=0}^{N-1}
      \frac{\rho^\beta-\rho^{\alpha+a\delta}}
           {1-\rho^{\alpha+a\delta}}.
\]
In particular, an even-power error expansion in the original step size uses
\(\alpha=\delta=2\), and hence the Gaussian base is \(\rho^2\),
not \(\rho\).

\section{Logarithmic errors and repeated spectral roots}
The geometric operator also handles polynomial--logarithmic errors. The key
identity is
\begin{equation}\label{qcc:log-spectral-identity}
 p(T_\rho)\bigl(h^\beta(\log h)^r\bigr)
  =\frac{\partial^r}{\partial\beta^r}
            \bigl(h^\beta p(\rho^\beta)\bigr),\qquad h>0.
\end{equation}
Indeed, \(h^\beta(\log h)^r=\partial_\beta^rh^\beta\), and
\(p(T_\rho)\) is a finite sum of substitutions independent of \(\beta\).

\begin{proposition}[Confluent geometric cancellation]
\label{qcc:confluent-filter}
If \(p(z)\) has a zero of multiplicity at least \(s\) at
\(z=\rho^\beta\), then \(p(T_\rho)\) annihilates
\(h^\beta(\log h)^r\) for every \(0\le r<s\).
For distinct positive exponents \(\beta_1,\ldots,\beta_A\) and positive
integers \(s_1,\ldots,s_A\), the normalized filter
\begin{equation}\label{qcc:confluent-polynomial}
 p(z)=\prod_{a=1}^{A}
        \left(\frac{z-\lambda_a}{1-\lambda_a}\right)^{s_a},
 \qquad\lambda_a=\rho^{\beta_a},
\end{equation}
preserves constants and cancels all the corresponding logarithmic modes.
Its coefficients have the recurrence-free Bell formula
\begin{align}
 C_0&=\prod_{a=1}^{A}
       \frac{(-\lambda_a)^{s_a}}{(1-\lambda_a)^{s_a}},
       \label{qcc:confluent-C}\\
 [z^\ell]p(z)&=\frac{C_0}{\ell!}
   \Y\ell\left(-0!\sum_a s_a\lambda_a^{-1},
                -1!\sum_a s_a\lambda_a^{-2},\ldots,
                -(\ell-1)!\sum_a s_a\lambda_a^{-\ell}\right).
       \label{qcc:confluent-Bell}
\end{align}
For \(\ell>\sum_a s_a\), the latter expression is identically zero.
\end{proposition}
\begin{proof}
The derivative of \(\rho^\beta\) with respect to \(\beta\) is
\((\log\rho)\rho^\beta\ne0\). A zero of order \(s\) in the
spectral coordinate therefore remains a zero of order \(s\) in the
exponent coordinate. Applying \eqref{qcc:log-spectral-identity} with
\(r<s\) gives zero. Formula \eqref{qcc:confluent-polynomial} has all
these zeros and satisfies \(p(1)=1\).

For its coefficients, write
\[
 p(z)=C_0\prod_{a=1}^{A}(1-z/\lambda_a)^{s_a},\qquad
 \log\frac{p(z)}{C_0}
   =-\sum_{r\ge1}\left(\sum_as_a\lambda_a^{-r}\right)\frac{z^r}{r}.
\]
Exponentiation is exactly the Bell master formula. The degree assertion follows
because \(p\) is a polynomial of degree \(\sum_as_a\).
\end{proof}

For an asymptotic expansion rather than a convergent series, the cancellation
claims apply to any finite collection of terms for which a remainder estimate
has been established. A finite filter transports that remainder by the triangle
inequality. One should not apply an infinite formal expansion as an equality of
functions without a separate convergence or asymptotic argument.

\section{Connection with the Fabius-function setting}
The relevant bridge to the surrounding repository is the geometric grid and its
spectral algebra. Whenever an approximation to a Fabius- or Rvachev-related
quantity is known to have an expansion in powers of a dyadic scale, its known
error terms can be canceled by the corresponding filter. For a grid ratio of
\(1/2\), integral powers use Gaussian base \(1/2\); an even-power
expansion uses base \(1/4\). Logarithmic corrections require repeated roots,
as just proved.

This application is deliberately conditional on the expansion of the particular
approximation being used. No assertion here identifies a specific Fourier
normalization, spline sequence, or exact dyadic Fabius value with such an error
series without proving that additional bridge. The unconditional result is the
filter algebra, including explicit weights, complete power response, analytic
remainder bound, and noise norm. That separation makes the result reusable
without importing hidden normalization assumptions.

\chapter{Multinomial coefficients and factorial ratios}
\label{qcc:factorial-ratios}
\section{The Gaussian multinomial extension}
For \(k_1+\cdots+k_s=n\), define
\[
 G_{k_1,\ldots,k_s}(q)
  =\frac{(q;q)_n}{\prod_{i=1}^{s}(q;q)_{k_i}}
  =\GB n{k_1}\GB{n-k_1}{k_2}\cdots
                  \GB{n-k_1-\cdots-k_{s-1}}{k_s}.
\]
The product expression proves polynomiality, nonnegative integer coefficients,
and the value \(n!/\prod_i k_i!\) at one. The degree is
\[
 D=\sum_{i<j}k_ik_j=\frac12\left(n^2-\sum_i k_i^2\right).
\]
The last expression follows by adding the degrees in the product; reciprocity
follows by multiplying the Gaussian reciprocity identities.

\begin{proposition}[Multinomial coefficient and parameter kernels]
\label{qcc:multinomial-kernels}
For \(r\ge1\), put
\[
 A_r=\sum_{i=1}^{s}\sum_{\substack{d\mid r\\d\le k_i}}d
           -\sum_{\substack{d\mid r\\d\le n}}d.
\]
Then
\begin{equation}\label{qcc:multinomial-coefficients}
 [q^R]G_{k_1,\ldots,k_s}(q)
   =\frac1{R!}\Y R\bigl(0!A_1,1!A_2,\ldots,(R-1)!A_R\bigr).
\end{equation}
Its logarithmic-parameter cumulants at one are
\begin{equation}\label{qcc:multinomial-cumulants}
 \kappa_1=\frac D2,\qquad
 \kappa_r=\frac{\beta_r}{r}
       \left(P_r(n)-\sum_{i=1}^{s}P_r(k_i)\right)\quad(r\ge2),
\end{equation}
where \(P_r(M)=\sum_{j=1}^{M}j^r\). Its multiplicity at a primitive
\(d\)-th root is
\begin{equation}\label{qcc:multinomial-carry}
 \nu_d=\left\lfloor\frac nd\right\rfloor
            -\sum_{i=1}^{s}\left\lfloor\frac{k_i}{d}\right\rfloor
          \in\{0,1,\ldots,s-1\}.
\end{equation}
\end{proposition}
\begin{proof}
For the coefficient formula, take the formal logarithm at zero. Each
\(\log(1-q^j)\) contributes \(-j/r\) to the coefficient of \(q^r\)
when \(j\mid r\), and zero otherwise. Subtracting denominator
logarithms gives \(A_r/r\), and Bell expansion proves
\eqref{qcc:multinomial-coefficients}. The calculation at one is the
factor-by-factor regular logarithm of \cref{qcc:regular-log}, with
\(P_r(n)-\sum_iP_r(k_i)\) in place of the two-block power-sum difference.
For the multiplicity, count multiples of \(d\) in the factorials.
Writing \(k_i=a_id+b_i\), \(0\le b_i<d\), gives
\(\nu_d=\lfloor\sum_i b_i/d\rfloor\), proving the stated range.
\end{proof}

Unlike a Gaussian binomial polynomial, a Gaussian multinomial need not be
squarefree. For example, the polynomial \(G_{1,1,1,1}(q)=[4]_q!\) has multiplicity
two at \(-1\):
\begin{equation}\label{qcc:multinomial-repeat-example}
 G_{1,1,1,1}(q)=(1+q)^2(1+q+q^2)(1+q^2).
\end{equation}
For its local jets, define \(C_\zeta\) and \(\lambda_r\) by taking
all exponents \(1,\ldots,n\) in the numerator and all exponents
\(1,\ldots,k_i\), with multiplicities, in the denominator of
\eqref{qcc:root-C} and \eqref{qcc:root-lambda}. Then
\eqref{qcc:root-factorization} and
\eqref{qcc:root-ordinary-derivatives} hold verbatim with
\(\nu=\nu_d\). Their proof was factorwise and never required
\(\nu\le1\). In this way the same Bell formula handles higher-order zeros
without a new inversion calculation.

\subsection{A multinomial Lucas formula}
Let \(n=ad+b\), \(k_i=a_id+b_i\), with all remainders in
\(\{0,\ldots,d-1\}\). If \(\sum_i b_i>b\), then \(\nu_d>0\)
and the multinomial vanishes at a primitive \(d\)-th root. If
\(\sum_i b_i=b\), there is no carry and
\begin{equation}\label{qcc:multinomial-lucas}
 G_{k_1,\ldots,k_s}(\zeta)
 =\binom{a}{a_1,\ldots,a_s}\,
                 G_{b_1,\ldots,b_s}(\zeta).
\end{equation}
Indeed, use the product of Gaussian binomials and apply \(q\)-Lucas at
each step. Without a carry, the ordinary binomial factors multiply to the
ordinary multinomial, and the remainder factors multiply to the smaller
Gaussian multinomial. With a carry, the multiplicity calculation already
proves vanishing.

\section{A polynomiality criterion for balanced factorial ratios}
Let \(a_1,\ldots,a_u\) and \(b_1,\ldots,b_v\) be nonnegative integers
with \(\sum_i a_i=\sum_j b_j\), and set
\[
 F(q)=\frac{\prod_{i=1}^{u}(q;q)_{a_i}}
                 {\prod_{j=1}^{v}(q;q)_{b_j}}.
\]
Let \(M\) be the largest index, taking \(M=0\) if every index is zero.

\begin{theorem}[Finite cyclotomic criterion]\label{qcc:ratio-criterion}
The rational function \(F(q)\) is a polynomial if and only if
\begin{equation}\label{qcc:ratio-floors}
 E_d:=\sum_i\left\lfloor\frac{a_i}{d}\right\rfloor
       -\sum_j\left\lfloor\frac{b_j}{d}\right\rfloor\ge0
       \qquad(2\le d\le M).
\end{equation}
When these conditions hold,
\begin{equation}\label{qcc:ratio-cyclotomic}
 F(q)=\prod_{d=2}^{M}\Phi_d(q)^{E_d}\in\Z[q],\qquad F(0)=1.
\end{equation}
Its degree and value at one are
\begin{equation}\label{qcc:ratio-degree-value}
 D=\frac12\left(\sum_i a_i^2-\sum_j b_j^2\right),\qquad
 F(1)=\frac{\prod_i a_i!}{\prod_j b_j!}.
\end{equation}
It satisfies \(F(q^{-1})=q^{-D}F(q)\). All its parameter jets follow
from the same factorwise deflation as \cref{qcc:root-jets}.
\end{theorem}
\begin{proof}
Cyclotomic factorization gives \(F=\prod_{d=1}^{M}\Phi_d^{E_d}\).
The balancing hypothesis makes \(E_1=0\) and cancels the signs from
\(1-q^a\). Distinct cyclotomic factors have disjoint roots. A negative
\(E_d\) would therefore give a genuine pole and prevent polynomiality;
nonnegative exponents give the integer polynomial
\eqref{qcc:ratio-cyclotomic}. The degree is the difference of the sums
\(1+\cdots+a_i\) and \(1+\cdots+b_j\); balancing cancels their linear
parts. Cancelling \(1-q\) from every individual factor before taking
\(q\to1\) gives the factorial value. Finally,
\(1-q^{-r}=-q^{-r}(1-q^r)\), so the same degree and sign calculation
gives reciprocity. The proof of \cref{qcc:root-jets} applies to this
multiset of factors without change.
\end{proof}

Polynomiality does not imply coefficientwise positivity. A small balanced
example is
\[
 \frac{(q;q)_6(q;q)_1^2}{(q;q)_5(q;q)_3}
 =\frac{(1-q^6)(1-q)}{(1-q^3)(1-q^2)}
 =1-q+q^2=\Phi_6(q).
\]
Thus the finite-field and partition interpretations of Gaussian coefficients
supply information beyond the cyclotomic criterion. Conversely, the criterion
extends the local-series machinery to examples for which no positive counting
interpretation is available.

\chapter{Computation, verification, and integration}
\label{qcc:computation}
\section{Choosing a representation for the task}
The formulae above are equal algebraically, but they serve different
computational purposes. A recurrence-free expression is a useful mathematical
specification; it is not automatically the fastest or best-conditioned numerical
algorithm.

For an exact Gaussian polynomial, the division-free Pascal recurrence is usually
a natural baseline: it uses only integer polynomial addition and shifts. For a
single coefficient in the base, \cref{qcc:gaussian-coefficients} supplies a
finite divisor-sum Bell expression or a restricted-partition expression. For a
coefficient of a product in an auxiliary variable, the geometric-product master
formula avoids expanding all factors separately. One first computes a finite
list of power sums and then evaluates one Bell polynomial.

For symbolic derivatives at a generic parameter, the logarithmic derivatives
\(L_r\) in \cref{qcc:generic-jets} compress large product-rule expressions.
At one, Bernoulli cumulants are simpler. At another root of unity, use the
constant, multiplicity, and regular logarithmic data \((C_\zeta,\nu,\lambda_r)\)
from \cref{qcc:root-jets}. Do not use a product quotient as a numerical
\(0/0\) and hope that floating-point cancellation will reconstruct the limit.

For basis conversion, the triangular \(q\)-Stirling recurrences give a
convenient algorithm, while the finite Gaussian sum and Bell diagonal formulae
give independent specifications. For implicit series, compute the logarithmic
power sums of the defining product and apply \cref{qcc:inverse-product}; the
ordinary Lagrange factor remains \(1/n\). For extrapolation, compute the
weights directly from \eqref{qcc:weights} and use the error and noise formulae
together, rather than increasing the cancellation order without controlling
sample accuracy.

\section{Arithmetic cost and finite verification}
Suppose the power sums \(P_1,\ldots,P_R\) are available. The coefficients
\(f_0,\ldots,f_R\) of
\(F(z)=\exp(\sum_{r\ge1}P_rz^r/r)\) can be evaluated through
\begin{equation}\label{qcc:coefficient-algorithm}
 f_0=1,\qquad nf_n=\sum_{r=1}^{n}P_rf_{n-r}.
\end{equation}
This uses \(O(R^2)\) arithmetic operations and \(O(R)\) stored scalars.
It follows by differentiating \(F\) and comparing coefficients in
\(zF'(z)=F(z)\sum_{r\ge1}P_rz^r\). The algorithm is recurrent, but
its output is specified independently by the recurrence-free Bell sum
\eqref{qcc:master-coefficient}. The distinction permits efficient calculation
and an independent formula audit at the same time. Bit complexity also depends
on coefficient heights and is not described by the arithmetic-operation count.

The companion \texttt{verify\_identities.py} uses only Python's standard
library, integer polynomial arithmetic, and exact rational numbers. Its successful
run for this article reports
\begin{center}
\fbox{\large\bfseries 10,912 exact checks in 36 groups.}
\end{center}
The tests include Gaussian products and convolution, both terminating
\(q\)-Chu--Vandermonde sums, divisor identities, Jackson stencils and Taylor
reconstruction, \(q\)-Stirling sums and Bell diagonals, weighted composition,
Lagrange coefficients, cyclotomic factors, parameter cumulants, all-root local
jets, and ordinary and logarithmic extrapolation filters.

The checks deliberately use different representations. Gaussian polynomials are
constructed by polynomial Pascal addition and compared with product values;
Bell expressions are checked against direct finite products or moments;
Lagrange coefficients are compared with ordinary truncated fixed-point
substitution; and \(q\)-Lucas is checked as exact divisibility by a
cyclotomic polynomial. Root-of-unity jets are checked not just at \(-1\),
but in the exact fields \(\Q(\zeta_d)\), \(3\le d\le8\), using
polynomials modulo \(\Phi_d\). For \(u\) of exact order \(L>1\),
the finite certificate
\[
 (1-u)^{-1}=-\frac1L\sum_{j=1}^{L-1}ju^j
\]
provides exact division in these calculations: multiplication by \(1-u\)
reduces the numerator sum to \(-L\).

Polynomial coefficient and moment tests cover \(0\le n\le12\), all
\(0\le k\le n\), and orders through eight. The \(-1\) tests cover
orders through nine. The cyclotomic-field jet tests cover \(0\le n\le9\),
all lower indices, and orders through seven. Other test bounds and the exact
rational parameter values are specified directly in the program. The archive
includes the complete output as \texttt{validation\_report.txt}.

\begin{auditbox}[title={What this verification establishes}]
These deterministic finite checks are reproducible audits of signs, shifts,
normalizations, and coefficient formulae. They are not proofs for unbounded
indices, do not establish floating-point stability outside the stated bounds,
and do not assert Lean verification. The mathematical proofs in the preceding
chapters establish the general statements. No current declaration inventory or
formalization status of the linked repository is inferred from these tests.
\end{auditbox}

\section{A theorem-level insertion map}
The following map identifies the mathematical interface with the source without
assuming that its chapter numbering will remain fixed.
\begin{longtable}{@{}p{0.31\textwidth}p{0.63\textwidth}@{}}
\toprule
Source mechanism & Extension supplied here\\
\midrule
\endfirsthead
\toprule
Source mechanism & Extension supplied here\\
\midrule
\endhead
Newton expansion and factorial bases &
Geometric Jackson--Newton interpolation, explicit Gaussian stencils, and the
\([j]_q\)-node \(q\)-Stirling arrays
(\cref{qcc:jackson,qcc:qstirling}).\\[2mm]
Exponential formula and Bell polynomials &
Finite-alphabet power sums, arbitrary geometric-product exponents, determinants,
and inverse logarithmic coordinates (\cref{qcc:bell-master}).\\[2mm]
Binomial inversion and incidence algebra &
Gaussian inversion, subspace-lattice M\"obius coefficients, and positive
linearization (\cref{qcc:convolution}).\\[2mm]
Stirling, Bernoulli, and Eulerian families &
Parameter derivatives and cumulants, ordinary/logarithmic derivative conversion,
and regular root-of-unity factors
(\cref{qcc:parameter-jets,qcc:roots}).\\[2mm]
Riordan and composition calculus &
A specified \(q\)-factorial normalization of ordinary composition and its
matrix group law (\cref{qcc:composition}).\\[2mm]
Lagrange--B\"urmann and Good inversion &
Product-defined inverses with finite Bell coefficients and a fully computed
cross-coupled two-variable case (\cref{qcc:lagrange}).\\[2mm]
Asymptotic coefficient operators &
Geometric spectral filters, complete error coefficients, noise norms, and
confluent cancellation of logarithmic modes (\cref{qcc:filters}).\\
\bottomrule
\end{longtable}

\section{Concluding perspective}
Gaussian coefficients occur here in three complementary roles. They are
coefficients of elementary and complete symmetric functions on geometric
alphabets; they are kernels for triangular basis changes and inversions; and
they are the explicit weights of spectral polynomials on geometric grids.
Ordinary Bell polynomials connect these roles by turning logarithmic power sums
into finite coefficients. Stirling numbers change the derivative coordinate,
Bernoulli numbers remove the apparent singularity at one, and Eulerian
polynomials make the remaining logarithmic derivatives rational.

This common mechanism is more informative than a formal replacement of every
integer by a \(q\)-integer. It explains which structures survive unchanged,
which acquire a Gaussian kernel, and which require a different operation.
It also separates exact finite algebra from local analysis and asymptotic
remainder estimates. Those separations are what allow the extension to serve
as a reusable part of the source's coefficient calculus.

\appendix
\chapter{A compact formula atlas}
\label{qcc:atlas}
This atlas records the input and output of the principal coefficient mechanisms.
All symbols and domains are defined in the referenced results.

\section*{Finite geometric products}
For
\(F(z)=\prod_\ell(a_\ell z;q^{d_\ell})_{N_\ell}^{-\alpha_\ell}\), put
\(P_r=\sum_\ell\alpha_\ell a_\ell^r[N_\ell]_{q^{d_\ell r}}\). Then
\[
 [z^n]F(z)=\frac1{n!}\Y n\bigl(((r-1)!P_r)_{r=1}^{n}\bigr)
 =\sum_{\sum rm_r=n}\prod_{r=1}^{n}\frac{P_r^{m_r}}{r^{m_r}m_r!}.
\]
This is \cref{qcc:master-theorem}. For an implicit solution \(y=zF(y)\),
replace \(P_r\) by \(nP_r\), use Bell degree \(n-1\), and divide by
\(n!\), as in \cref{qcc:inverse-product}.

\section*{Gaussian polynomial coefficients}
For \(m=n-k\), set
\[
 A_r=\sum_{\substack{d\mid r\\1\le d\le k}}d
     -\sum_{\substack{d\mid r\\m<d\le n}}d.
\]
Then
\[
 [q^R]\GB nk=\frac1{R!}
       \Y R\bigl(((r-1)!A_r)_{r=1}^{R}\bigr).
\]
The two-block divisor kernels are the multiset difference of denominator and
numerator exponents; see \cref{qcc:gaussian-coefficients}.

\section*{The \(q=1\) expansion}
Let
\(\kappa_1=k(n-k)/2\) and
\(\kappa_r=(\beta_r/r)\sum_{j=1}^{k}((n-k+j)^r-j^r)\), \(r\ge2\).
Then
\[
 \GB[e^t]nk=\binom nk\sum_{r\ge0}\Y r(\kappa_1,\ldots,\kappa_r)
                                      \frac{t^r}{r!}.
\]
For ordinary \(q\)-derivatives, apply signed Stirling coefficients as in
\eqref{qcc:ordinary-jets-one}.

\section*{All roots of unity}
Deflate every factor \(1-\zeta^ae^{at}\), using
\(-at\) when \(\zeta^a=1\) and \(1-\zeta^a\) otherwise. The constant
product is \(C_\zeta\), the net order is \(\nu\), and the regular
logarithmic derivatives are \(\lambda_r\). Then
\[
 \GB[\zeta e^t]nk=C_\zeta t^\nu
      \sum_{r\ge0}\Y r(\lambda_1,\ldots,\lambda_r)\frac{t^r}{r!}.
\]
Equations \eqref{qcc:eta}--\eqref{qcc:root-lambda} supply every input by
finite sums, and \eqref{qcc:root-ordinary-derivatives} converts to ordinary
Taylor coefficients.

\section*{Geometric extrapolation}
For \(0<\rho<1\), the weights
\[
 w_{N,j}=\frac{(-1)^{N-j}\rho^{\binom{N-j+1}{2}}
                  \GB[\rho]Nj}{(\rho;\rho)_N}
\]
preserve constants, kill the first \(N\) powers, and have exact absolute norm
\(( -\rho;\rho)_N/(\rho;\rho)_N\). Their response to every higher power
is \((-1)^N\rho^{\binom{N+1}{2}}\GB[\rho]{m-1}{N}\).
Repeated spectral roots cancel logarithmic modes; see
\cref{qcc:confluent-filter}.

\chapter{Source provenance and reproducibility}
\label{qcc:provenance}
\section{Source and scope}
The requested source is the canonical \emph{Combinatorial Coefficient Calculus}
under the repository path given in \cite{source}. The directory metadata read
for this extension identified its \LaTeX{} file by Git blob SHA
\begin{center}\small
\texttt{ffbe7c117db3fbeb3c0ae804e84422bb2dc6a89b}.
\end{center}
This is a content identifier for that file, not a claim about the repository's
latest commit or its formalization coverage. The source interface used here is
its coefficient rules, Newton basis, Bell generating functions, Stirling
normalizations, and composition/inversion organization. The present document is
not a line-by-line revision of the source, and the original files were not
modified.

The classical identities cited in the bibliography are proved in the text in
the normalizations actually used. The root-jet and filter developments are
presented with complete derivations, rather than a priority claim. No inference
of novelty is made from the absence of a result in a small set of references.

\section{Files and commands}
The archive contains the standalone \LaTeX{} source, its compiled PDF,
\texttt{verify\_identities.py}, \texttt{validation\_report.txt}, and a
\texttt{README.md}. No bibliography processor or repository-specific notation
file is needed. With a \TeX{} installation containing Libertinus and the listed
standard packages, compile with
\begin{verbatim}
pdflatex -interaction=nonstopmode -halt-on-error \
  q_binomial_coefficient_calculus.tex
pdflatex -interaction=nonstopmode -halt-on-error \
  q_binomial_coefficient_calculus.tex
\end{verbatim}
A third pass may be needed after substantial edits to the table of contents.
Equivalently, use \texttt{latexmk -pdf q\_binomial\_coefficient\_calculus.tex}.
The verifier requires Python 3.9 or later and no third-party packages:
\begin{verbatim}
python verify_identities.py
\end{verbatim}
The supplied report records the successful run accompanying this PDF. Running
the program again recomputes the certificates; it does not read or trust the
saved report.

\backmatter
\begin{thebibliography}{9}
\addcontentsline{toc}{chapter}{Bibliography}

\bibitem{source}
Vladimir Reshetnikov, \emph{ProveIt repository: Combinatorial Coefficient
Calculus}, canonical \LaTeX{} monograph and accompanying provenance files.
Source consulted September 4, 2026.
\url{https://github.com/VladimirReshetnikov/ProveIt/tree/main/Analysis/FabiusFunction/docs/semi-formalized-research-frontiers/drafts/combinatorial-coefficient-calculus/Combinatorial_Coefficient_Calculus}.

\bibitem{dlmf172}
NIST Digital Library of Mathematical Functions, \S17.2,
\emph{q-Calculus}. Definitions of shifted factorials, Gaussian coefficients,
finite binomial identities, and Jackson differentiation.
\url{https://dlmf.nist.gov/17.2}.

\bibitem{dlmf175}
NIST Digital Library of Mathematical Functions, \S17.5,
\emph{\({}_0\phi_0,{}_1\phi_0,{}_1\phi_1\) Functions}.
Infinite \(q\)-binomial identities and their conventions.
\url{https://dlmf.nist.gov/17.5}.

\bibitem{dlmf176}
NIST Digital Library of Mathematical Functions, \S17.6,
\emph{\({}_2\phi_1\) Function}. Terminating \(q\)-Chu--Vandermonde
summations.
\url{https://dlmf.nist.gov/17.6}.

\bibitem{cai}
Yue Cai, Richard Ehrenborg, and Margaret Readdy,
``\(q\)-Stirling identities revisited,''
\emph{Electronic Journal of Combinatorics} \textbf{25}(1) (2018), Paper P1.37.
\url{https://doi.org/10.37236/6829}.
Preprint: \url{https://arxiv.org/abs/1706.06623}.

\bibitem{dilcher}
Karl Dilcher, ``Some \(q\)-series identities related to divisor functions,''
\emph{Discrete Mathematics} \textbf{145} (1995), 83--93.
\url{https://www.sciencedirect.com/science/article/pii/0012365X9500092B}.

\bibitem{formichella-straub}
Sam Formichella and Armin Straub,
``Gaussian binomial coefficients with negative arguments,''
\emph{Annals of Combinatorics} \textbf{23}(3) (2019), 725--748.
\url{https://arxiv.org/abs/1802.02684}.
Author publication record:
\url{https://arminstraub.com/pub/qbinomial-negative}.

\bibitem{gessel}
Ira M. Gessel, ``Lagrange inversion,''
\emph{Journal of Combinatorial Theory, Series A} \textbf{144} (2016), 212--249.
\url{https://arxiv.org/abs/1609.05988}.

\end{thebibliography}
\end{document}
